\documentclass[12pt]{article}

%Article standard packages
\usepackage[latin9]{inputenc}
\usepackage{mathpazo}
\usepackage{geometry}
\usepackage{caption,subfig}
\usepackage{mathrsfs}
\usepackage{parskip}
\geometry{verbose}
\usepackage[colorlinks,citecolor={black},linkcolor={black},urlcolor={green}]{hyperref}
\usepackage{color}
\usepackage{amsmath}
\usepackage{amssymb}
\usepackage{graphicx}
\usepackage{afterpage}
\usepackage{float}
\usepackage[section]{placeins}
\usepackage{pdfpages}
\usepackage{tikz} 
\usepackage{fancyhdr}
\usepackage[labelfont=bf]{caption}
\usepackage{color}
\PassOptionsToPackage{normalem}{ulem}
\usepackage{ulem}
\usepackage{dsfont}
\usepackage{threeparttable}
\usepackage{enumitem}  
\usepackage{tcolorbox}
\usepackage{multicol}
\usepackage{tabularx}
\usepackage{booktabs}
\usepackage{lscape}
\usepackage{float}
\usepackage{caption}
\usepackage{listings}
\usepackage{booktabs}
\usepackage{xcolor}
\usepackage{multicol}
\usepackage{dsfont}
\usepackage{color}
\usepackage{eurosym}
\usepackage{amsfonts}
\usepackage{rotating}
\usepackage{lscape}
\usepackage{comment}
\usepackage{natbib}
\usepackage{setspace}
\usepackage{pgfplots}
\usepackage{multirow}
\usepackage[linesnumbered]{algorithm2e}
\usepackage{amsthm}     
\usepackage{bm}         

\makeatletter
\newcommand{\multiline}[1]{%
  \begin{tabularx}{\dimexpr\linewidth-\ALG@thistlm}[t]{@{}X@{}}
    #1
  \end{tabularx}
}
\makeatother
\makeatletter
\providecommand{\tabularnewline}{\\}
\usepackage{afterpage}
%All other maths elements packages
\setcounter{MaxMatrixCols}{10}
\newtheorem{theorem}{Theorem}
\newtheorem{acknowledgement}{Acknowledgement}
\newtheorem{axiom}{Axiom}
\newtheorem{case}{Case}
\newtheorem{ass}{Assumption}
\newtheorem{claim}{Claim}
\newtheorem{conclusion}{Conclusion}
\newtheorem{condition}{Condition}
\newtheorem{conjecture}{Conjecture}
\newtheorem{corollary}{Corollary}
\newtheorem{criterion}{Criterion}
\newtheorem{definition}{Definition}
\newtheorem{example}{Example}
\newtheorem{exercise}{Exercise}
\newtheorem{lemma}{Lemma}
\newtheorem{notation}{Notation}
\newtheorem{problem}{Problem}
\newtheorem{proposition}{Proposition}
\newtheorem{remark}{Remark}
\newtheorem{solution}{Solution}
\newtheorem{summary}{Summary}
\geometry{left=1in,right=1in,top=1in,bottom=1in}
\newdimen\dummy
\dummy=\oddsidemargin
\addtolength{\dummy}{72pt}
\marginparwidth=.5\dummy
\marginparsep=.1\dummy
\marginparsep=.1\dummy
\makeatother
\usepackage{amsmath}
\DeclareMathOperator*{\argmax}{arg\,max}
\DeclareMathOperator*{\argmin}{arg\,min}
\DeclareMathOperator{\Gr}{Gr}
\newcommand{\Break}{\State \textbf{break} }

\newsavebox\CBox
\newcommand\hcancel[2][0.5pt]{%
  \ifmmode\sbox\CBox{$#2$}\else\sbox\CBox{#2}\fi%
  \makebox[0pt][l]{\usebox\CBox}%  
  \rule[0.5\ht\CBox-#1/2]{\wd\CBox}{#1}}
  
\makeatletter
\newcommand{\proofpart}[2]{%
  \par
  \addvspace{\medskipamount}%
  \noindent\emph{Part #1: #2}\par\nobreak
  \addvspace{\smallskipamount}% 
  \@afterheading
}
\makeatother

\usepackage[USenglish]{isodate}% http://ctan.org/pkg/isodate

%% Notation 

\newcommand{\DChSet}{\mathbb{D}}
\newcommand{\dVecVar}{\mathbf{d}}
\newcommand{\mVecVar}{\mathbf{m}}
\newcommand{\lVecVar}{\mathbf{l}}
\newcommand{\qVecVar}{\mathbf{q}}
\newcommand{\sVecVar}{\mathbf{s}}
\newcommand{\rVecVar}{\mathbf{r}}


\begin{document}

%%%%%%%%%%%%%%%%%%%%%%%%%%%%%%%%%%%%%%%%%%%%%

\title{\vspace{-1cm} A fast upper envelope scan for \\ discrete-continuous dynamic programming\thanks{\footnotesize{}{}{}
We thank Hazel Bateman, Chris Carroll, Mike Keane, Ben Newell, Ben Moll, John Rust, Fedor Ishkakov, Greg Kaplan, Bertel Schjerning, Susan Thorp, Alan Lujan, Anant Mathur, Adrian Monninger, Mateo Velasquez-Giraldo, Matt White and conference and seminar participants at several institutions for valuable feedback, discussions, and comments. Research support from the Australian Research Council (ARC LP150100608) and the ARC Center of Excellence in Population Aging Research (CE17010005) is appreciated. This research was undertaken with the assistance of resources and services from the National Computational Infrastructure (NCI), which is supported by the Australian Government. The software implementing FUES in our applications, as well as the one providing the function that implements FUES in its general form can be found at \href{http://github.com/akshayshanker/FUES_EGM}{\textcolor{blue}{here}}.}}
\author{Loretti I. Dobrescu\thanks{{\footnotesize{}{}{}Corresponding author. School of Economics, University of New South Wales, Australia. Email: dobrescu@unsw.edu.au.}} \, and \ Akshay Shanker \thanks{{School of Economics, University of New South Wales, Australia. Email: a.shanker@unsw.edu.au.}}
}
\date{} %<--- this removes the automatic date 
%
\maketitle

\begin{center}
First draft: August 4, 2022 \par
Current draft: \today
\vspace{3em}
\end{center}

\begin{abstract}
We develop a general fast upper envelope scan algorithm to solve any stochastic dynamic models with discrete-continuous choices where the optimization problem has a one-dimensional representation. For this vast class of models, FUES enables the endogenous grid method by efficiently selecting the optimal solution using a targeted secant scan. Importantly, it does so without requiring monotonicity of the policy function. Using several workhorse applications in the literature, we show the massive computational speed and accuracy gains generated across the board. Importantly, FUES advances the frontier of problems that can be handled by Euler-equation techniques: when applied to a discrete housing model with present-bias individuals, FUES uncovers a staircase-shaped saving rule with locally zero marginal propensities to save. Finally, we prove that FUES accurately recovers the approximate upper envelope in the general case under standard economic conditions. 


\vspace{1em}
\textbf{Key Words}: discrete-continuous choices, non-convex optimization, Euler equation, stochastic dynamic programming

\textbf{JEL Classification}: C13, C63, D91
\end{abstract}

\newpage

\onehalfspacing

%%%%%%%%%%%%%%%%%%%%%%%%%%%%%%%%%%%%%%%%%%%%%

\begin{comment}
    -----------
    Terminology 
    ------------
    - upper envelope 
    - one-dimensional 
    - financial assets a_t
    - housing assets H_t
    - Total wealth (this is different to a_t and H_t!)
    - sequence of future discrete chocies
    - DC-EGM
    - choice-specific
    - life-cycle 
    - jump detection
    - Lipshitz
    - Application 1: Retirement choice
    - Application 2: Continuous housing investment with (adjustment) frictions (ok to use illiquid as an adjective). 
    - Application 3: Discrete housing choice under present-bias 
    - wage
    - Crossing point (not intersection point!). 
\end{comment}

\begin{comment}
---------------------
STATUS 23 June 2025
--------------------

Akshay penultiamte pass
-----------------------
[x] Abstract
[x] Intro

[x] Application 1
[x] Application 1 Numerical 
[x] Application 2
[ ] Application 2 Numerical 
[ ] Application 3
[ ] Application 3 Numerical 
[ ] Theory main section 
[ ] Theory Appendix
[ ] Tables, figures and numbers therein.
[ ] Refs to appendix and formatting. 

Bits and bobs
-------------
[x] Update Moll/Laibson/Maxted reference. 

\end{comment}

\section{Introduction}

Stochastic dynamic decision-making applications have increasingly been incorporating discrete-continuous choices. Such choices are central to understanding the dynamics of investment frictions \citep{Skiba1978, Khan2008}, housing adjustments \citep{Yogo2016a, Kaplan2020}, the implications of durable goods consumption for macroeconomic dynamics \citep{Kaplan2014,Berger2015}, life-cycle asset allocation \citep{Fagereng2019,Laibson2020}, labour supply \citep{Attanasio2018,ISKHAKOV2021401} or default risk \citep{Arellano2008,Arellano2016,Jang2023}, to name just a few areas of active research.

Currently, the standard approach to solve stochastic dynamic problems is the endogenous grid method (EGM).  Unlike value function iteration (VFI), EGM relies on analytically inverting the Euler equation \citep{Carroll2006} and may completely remove costly numerical root-finding or optimization steps \citep{Iskhakov2015}. The drawback however is that, similar to other Euler equation methods \citep{Coleman1990, Li2014}, EGM relies on the value function concavity to ensure Euler equation solutions are globally optimal. This concavity breaks down in models involving both continuous and discrete choices, leading to sub-optimal Euler equation solutions due to discontinuous policy functions \citep{Fella2014,Iskhakov2017}. Several methods have been suggested so far to calculate the upper envelope of the solutions generated by discrete-continuous EGM (DC-EGM; \cite{Fella2014, Iskhakov2017, Druedahl2017, Druedahl2021}).\footnote{We refer to the problem class as `DC-EGM' and reserve method names for solution algorithms.} All these methods, however, are complex and scale poorly, depend on monotonicity\footnote{Due to the monotonicity assumption, the benefits of EGM are unavailable whenever discrete-continuous choices create both upward and downward jumps -- or flat segments -- in the policy function, as in models with both durable and non-durable consumption \citep{Kaplan2014, Yogo2016a}, present-bias \citep{laibson21997}, or non-linear taxation \citep{Best2017, Gelber2020}.} of the policy function \citep{Iskhakov2017, Fella2014, Druedahl2017}, or rely on combining numerical optimization with Euler equation methods \citep{Druedahl2021} that can lead to loss of computational accuracy and efficiency. None supplies error bounds or a formal guarantee of recovering the true upper envelope. 

We address these challenges by developing a general and highly efficient method to compute the upper envelope of the value correspondence generated by EGM for problems where current period states can be collapsed to one-dimensional ones. Our fast upper envelope scan (FUES) method works by noting that the upper envelope of the value correspondence is the supremum of choice-specific concave value functions, with each value function corresponding to a history-dependent future sequence of discrete choices. The convex regions of the upper envelope occur where different choice-specific value functions cross, with the optimal policy function experiencing discontinuous jumps only in these regions. Accordingly, FUES performs a secant scan to determine whether including a potential optimal point forms a concave or convex region of the upper envelope. If it forms a concave region, the point is eliminated as sub-optimal if it induces a jump in the policy function. If, however, it forms a convex region of the upper envelope, the candidate point is considered optimal and retained.
 
We illustrate the wide applicability of our method via three applications. First, we implement FUES in a classic discrete-continuous dynamic model - i.e., a life-cycle model with discrete retirement and continuous consumption decisions \`a la \cite{Iskhakov2017}. Second, we implement FUES in a model of continuous housing investment with adjustment frictions where previous methods fail due to the non-monotonicity of the policy function. Featuring both liquid and illiquid assets, this example also shows how our method can be applied to EGM problems without an analytical inverse \citep{Iskhakov2015}. Finally, we implement FUES in a model with discrete housing and continuous financial assets \citep{Fella2014} in the presence of present-biased \citep{laibson21997} individuals -- a non-monotone application that has not previously been solved using any Euler equation method. 

For each application, we compare the performance of FUES the standard existing methods including the monotone segment selection (MSS) method proposed by \citep{Iskhakov2017}, the local triangulation method (LTM) proposed by \cite{Druedahl2021} and nested-EGM (NEGM, \cite{Druedahl2021}), as well as against the multidimensional rooftop-cut method (RFC, \cite{DS2024RFC}). In all cases, FUES demonstrates a considerable computational speed advantage. Where MSS and LTM can be applied, FUES accelerates total solution time by up to \textcolor{orange}{an order of magnitude} compared to MSS and more than \textcolor{orange}{two orders of magnitudes} compared LTM for larger grids. For instance, in a version of the discrete housing choice model without present bias (our third application), the computation time for the owner's problem takes \textcolor{orange}{1.1 seconds} using FUES, \textcolor{orange}{12.4 seconds using MSS} and over \textcolor{orange}{123 seconds using MSS}. We also find that computation time under FUES scales sub-linearly in grid size, while DC-EGM scales linearly and LTM scales quadratically. In the example with a continuous non-durable housing stock, LTM and MSS cannot be applied and compared with NEGM, we find numerically inverting the Euler equation and using FUES leads to an acceleration of about 50\%. In the discrete housing model with present-bias, we compare FUES to standard VFI as no existing Euler equation methods can be applied, and we find two orders of magnitude improvement in Euler error and speed. Across the board, speed advantages occur due to the lower complexity of FUES across model parameterizations, and without any parallelization or GPU acceleration. 

As for accuracy, FUES performs on average better than all existing methods. For each application, we compare the performance of each solution method against an \textcolor{orange}{ultra-fine grid search}. For monotone problems where an upper envelope method such as MSS and LTM can be applied, the accuracy gains of FUES are modest because both use the analytical inverse of the Euler equation. In non-monotone problems (Section \ref{sec: app2_cont_hous} and Section \ref{sec: Fella}),  FUES achieves over an order of magnitude improvement in accuracy compared to NEGM and two orders of magnitude improvement compared to VFI. 

Beyond computational gains, application of FUES to study the interaction of present-bias with housing decisions yields new economic insights into how financial savings and housing interact.  We show how a generalized Euler inequality in this environment need not hold with equality even when financial savings are strictly positive. In this case, we uncover regions in which sophisticated present-biased households deliberately hold a fixed amount of financial wealth as a commitment device to induce future housing investment. An implication is that local marginal propensities to consume are locally zero -- the individual holds just enough wealth to prevent future down-sizing by their impatient selves but does not wish to save any extra. 

Finally, we provide a theoretical foundation for FUES in its general form, showing how it can recover the upper envelope of the approximate value functions that form the EGM value correspondence without error when (i) the jumps between policy functions are bounded below (as is the case when there are finitely many jumps in the policy functions), (ii) the grid size is sufficiently large,\textcolor{orange}{and iii) the smooth regions of the problem satisfy standard concavity and Lipshitz conditions.} These conditions are straightforward to verify in practice, and ensure FUES can differentiate between jumps in the policy function that reflect shifts in future discrete choices and smooth changes in policy along a given choice sequence.

FUES advances considerably the latest methods to solve discrete-continuous dynamic models \citep{Fella2014, Iskhakov2017, Druedahl2021, DS2024RFC}. Both \cite{Fella2014} and  \cite{Iskhakov2017} rely on the monotonicity of the optimal policy functions, while \cite{Druedahl2021} requires a nested numerical optimization step. In contrast, the RFC implemented in \cite{DS2024RFC} requires analytical information on the gradient of the policy function and also performs a potentially complex nearest neighbor search at each grid point. In this respect, our method advances the literature in several ways. First, FUES is significantly less computational complex complex \citep{Fella2014, Iskhakov2017, Druedahl2021, DS2024RFC} than all alternative methods, offering a massive acceleration in computation time. Since FUES operates on (and returns) a single dimensional array regardless of model complexity, it is compatible with cutting-edge massively parallel and efficiently vectorized computing using languages such as \texttt{Google JAX}, \texttt{Numba} and \texttt{NumPy}. Second, FUES does not require the monotonicity of any of the policy functions, enabling its broad application to all discrete-continuous problems where the state-space is one dimensional. Third, unlike existing methods, FUES provides a theoretical guarantee of accuracy in the general case. Finally, FUES is considerably easier more intuitive to implement than existing upper envelope methods, and does not require knowing the shape of the policy functions or the gradient of the value function. Taken together, our numerical results suggest FUES is well-suited as a drop-in algorithm to compute the upper envelope in discrete-continuous problems featuring uncertainty, which covers a broad set problems found in recent literature, including those studying fiscal policy transmission \citep{brumm2024global}, the effect of durable goods on the marginal propensity to spend \citep{Kovacs2021,NBERw32080}, the distribution of homeownership \citep{LeBlanc2023housing} and labour market choices \citep{CRUCES2024104960, Bacher2024}. 

The paper proceeds as follows. In Section 2, we start by solving and discussing our main results in the context of three well-known applications. Section 2.1 introduces FUES informally using the simple retirement choice model in \cite{Iskhakov2017} where agents choose their savings and labor force participation; here we also compare the performance of FUES to the upper envelope methods proposed by \cite{Iskhakov2017} and \cite{Druedahl2017}.  Section 2.2 further demonstrates FUES in an application where existing upper envelope methods fail, using a model in which agents choose whether or not to adjust their illiquid housing assets. In Section 2.3, we further showcase the ability of FUES to handle non-monotone policies by solving a present-biased individual's problem in the discrete housing choice environment proposed by \cite{Fella2014}. In Section 3, we follow our study of these applications by formally stating FUES in its general form and providing the proofs on how \textcolor{orange}{FUES obtains the upper envelope}. Section 4 concludes.

%%%%%%%%%%%%%%%%%%%%%%%%%%%%%%%%%%%%%%%%%%%%%

\section{Illustrative applications}\label{sec: Applications}

We start by introducing FUES in the context of three well-known dynamic optimization applications. After briefly introducing each problem, we discuss how their discrete choices result in non-convexity, show how FUES can retrieve the optimal solution and compare FUES to existing methods. 

%%%%%%%%%%%%%%%%%%%

\subsection{Application 1: Retirement choice}\label{sec: retchoice}

\subsubsection{Model environment}

Consider the finite horizon retirement and savings choice model in \cite{Iskhakov2017} where an agent consumes, works (if they so choose), and saves from time $t=0$ until time $t=T$. At the beginning of each period, the agent starts as a worker or retiree, with the state variable denoting their beginning-of-period work status given by the discrete variable $d_{t}$. If the agent works, they earn a per-period wage $y$. Every period, the agent can choose to continue working during the next period by setting $d_{t+1} = 1$, or to permanently exit the workforce by setting $d_{t+1}=0$. If the agent chooses to work the next period, they will incur a utility cost $\tau$ at time $t$. We assume all agents start as workers so $d_{0} = 1$. Agents can also consume $c_{t}$ and save in financial assets $a_{t+1}$ (the continuation state), with $a_{t+1}\in S$ and $S\colon = \left[0,\bar{a}\right]\subset \mathbb{R}_{+}$.\footnote{\textcolor{orange}{Footnote here explaining use of terms state and continuation state and link to DCDP and Carroll 2025.}}  Thus, the intertemporal budget constraint is
%
\begin{equation}\label{eq: re_BC}
    a_{t+1} = (1+r)a_{t} + d_{t}y - c_{t}.
\end{equation}
%
Per-period utility is given by $\log(c_{t}) - \tau d_{t}$. Letting the function $u$ be defined by $u(c) = \log(c)$, the agent's maximization problem becomes
%
\begin{equation}
    V_{0}^{d_{0}}(a_{0}) = \max_{(c_{t},d_{t+1})_{t=0}^{T}}\left\{\sum_{t=0}^{T}\left[\beta^{t}u(c_{t}) - \tau d_{t+1}\right]\right\},
\end{equation}
%
subject to Equation \eqref{eq: re_BC}, $a_{t}\in S$ for each $t$ (the time $t$ or current state), and the fact that the agent cannot return to work after retiring (i.e., $d_{t+1} = 0$ if $d_{t} = 0$). Let $V_{t}^{d_{t}}$ denote the beginning-of-period value function. If the agent enters the period as a worker, the agent's time $t$ value function for any $t<T$ will be characterised by the Bellman equation
%
\begin{equation}\label{eq: bell1}
    V_{t}^{1}(a) = \max_{c,d_{t+1}\in \{0,1\}} \left\{u(c) - d_{t+1}\tau + \beta V_{t+1}^{d_{t+1}}(a^{\prime})\right\},
\end{equation}
%
where $a^{\prime} = (1+r)a + y - c$ and $a^{\prime}\in S$. For period $T$, the value function for the worker is $V_{T}^{1}(a) = u((1+r)a + y)$. If the agent enters the period as a retiree, the value function becomes
%
\begin{equation}
    V_{t}^{0}(a) = \max_{c} \left\{u(c) + \beta V_{t+1}^{0}(a^{\prime})\right\},
\end{equation}
%
with $a^{\prime} = (1+r)a - c$. For period $T$, the value function for the retiree is $V_{T}^{0}(a) = u((1+r)a)$. The optimization problem for the retiree is a standard concave problem. For the worker, however, the optimization problem is not concave since she optimizes jointly a discrete choice and a continuous choice. Moreover, even conditional on $d_{t+1} = 1$, next period value function $V_{t+1}^{1}$ will not be concave since the value function represents the supremum over \textit{all future feasible combinations of discrete choices}. The non-concavity of $V_{t+1}^{1}$ produces the `secondary kinks' described by \cite{Iskhakov2017}. 

To see how the choice at time $t$ implicitly controls the future sequence of discrete choices and produces the secondary kinks, write the time $t$ worker's value function as
%
\begin{equation}\label{eq: Bell_Q}
        V_{t}^{1}(a) = \max_{c}\max_{\textbf{d}\in \DChSet} \left\{u(c) + \beta Q_{t+1}^{\textbf{d}}(a^{\prime})\right\},
\end{equation}
%
where $Q_{t+1}^{\textbf{d}}$ is the $t+1$ value function conditional on a given sequence of future discrete choices $\textbf{d}$, with $\textbf{d} = \{d_{t+1}, d_{t+2},\dots,d_{T}\}$. In particular, letting $\boldsymbol{\bar{\tau}} = \{\tau, \beta\tau,\beta^{2}\tau,\dots,\beta^{T-1}\tau\}$, we have
%
\begin{equation}
    Q_{t+1}^{\textbf{d}}(a) = \max_{(c_{k})_{k=t+1}^{T}}\left\{\sum_{k=t+1}^{T}\beta^{k-t-1}u(c_{k})\right\} - \boldsymbol{\bar{\tau}}^{\text{tr}}\textbf{d}.
\end{equation}
%
\begin{figure}[htb!]
    \centering
    \captionsetup{width=.7\linewidth}
    \input{figures/concave_env.tikz}
    \caption{The time $t$ worker value function $V_{t}^{1}$ is the upper envelope of concave functions, where each concave function is a value function conditional on a \textit{sequence} of future discrete choices. The subscript `$i$' on $\mathbf{d}_{i}$ indicate distinct sequences of future discrete choices.}
    \label{fig: concav_env}
  \end{figure}

To sum up, the value function non-concavity, even holding the choice $d_{t+1}$ fixed, is brought on by the implicit changes in the entire future sequence of discrete choices as one controls the choice variable $c$ in \eqref{eq: Bell_Q}. In this case, the Bellman equation \eqref{eq: bell1} still holds, and one can numerically implement VFI to compute a solution. The challenge arises when solving the Bellman equation using numerical methods becomes burdensome computationally. A more computationally efficient strategy involves recovering the policy function by solving for points that satisfy the first order conditions (FOCs - i.e., the Euler equations) of the Bellman equation. However, since the upper envelope is not concave, the points satisfying the FOCs could be associated with any future sequence of discrete choice in Figure \ref{fig: concav_env}, and may not be on the upper envelope.

%%%%%%%%%%%%%%%%%%%

\subsubsection{The Euler equations}

We now discuss the FOCs and then proceed to our contribution (i.e., the FUES method) as a way to use necessary first order information to efficiently compute $V_{t}$. If the agent chooses $d_{t+1} =1$ (i.e., they continue as a worker in $t+1$), we can write the time $t$ worker Euler equation as
%
\begin{equation*}
    u^{\prime}(c_{t}^{1}) \geq  \beta (1+r) u^{\prime}(c_{t+1}),
\end{equation*}
%
where $c_{t}^{1}$ is the time $t$ consumption policy conditional on $d_{t+1} =1$, while $c_{t+1}$ is the unconditional time $t+1$ consumption policy. On the other hand, if the agent chooses $d_{t+1} = 0$ (i.e., they retire), then the Euler equation is
%
\begin{equation*}
    u^{\prime}(c_{t}^{0}) \geq \beta (1+r) u^{\prime}(c_{t+1}^{0}), 
\end{equation*}
where $c_{t}^{0}$ and $c_{t+1}^{0}$ are the consumption policies conditional on $d_{t+1} =0$ and $d_{t+2} = 0 $ respectively. 
% Clear above that retiree cannot work again?
%
\paragraph{Functional Euler equations.} It will now be helpful to write the Euler equation in its functional form. Let $\sigma_{t}^{d}\colon S\times \{0,1\}\rightarrow \mathbb{R}_{+}$ be the conditional continuation state policy function for the worker at time $t$ if $d=1$, and for the retiree if $d=0$. Note that $\sigma_{t}^{d}$ will depend, through its second argument, on the discrete choice (to work or not to work in $t+1$) made by the worker at time $t$. Unless required, for brevity, we refer to $\sigma_{t}^{d}$ simply as policy functions, with the `super-script' marking that it is a conditional on a current discrete choice. The time $t$ and time $t+1$ policy functions will satisfy the following functional Euler equation:
%
\begin{multline}\label{eq: Euler_func_ret_1}
    u^{\prime}((1+r)a + dy - \sigma^{d}_{t}(a,d_{t+1}))\\  \geq \beta (1+r) u^{\prime}((1+r)\sigma_{t}^{d}(a,d_{t+1}) + d_{t+1}y - \sigma_{t+1}^{d_{t+1}}(a^{\prime},d_{t+2})),
\end{multline}
%
where $a^{\prime}= \sigma_{t}^{d}(a,d_{t+1})$. 
%
On the choice of whether to work or not, the time $t$ worker will chose $d_{t+1} = 1$ if and only if
%
\begin{multline}\label{eq: disc_choice_t}
    u((1+r)a + y - \sigma_{t}^{1}(a, 1))  - \tau + \beta V_{t+1}^{1}(\sigma_{t}^{1}(a,1)) \\ > u((1+r)a + y- \sigma_{t}^{1}(a,0)) +  \beta V_{t+1}^{0}(\sigma_{t}^{1}(a,0)).
\end{multline}

Since the discrete choice is itself a function of the state, we can also define a discrete choice policy function $\mathcal{I}_{t}\colon S\times \{0,1\}\rightarrow \{0,1\}$. As such, we will have $d_{t+1} = \mathcal{I}_{t}(a,d)$ and $d_{t+2} = \mathcal{I}_{t+1}(a^{\prime},d_{t+1})$, where $\mathcal{I}_{t}$ is evaluated to satisfy \eqref{eq: disc_choice_t} each period conditional on the $t+1$ value function.  

%%%%%%%%%%%%%%%%%%%

\subsubsection{Computation using EGM and FUES}\label{sec: FUES_RET}

We now turn to how FUES can identify the upper envelope from a set of points that satisfy the Euler equations. Fix a time $t$ and suppose the value function $V_{t+1}^{d}$, the discrete choice function $\mathcal{I}_{t+1}$, and the optimal policy function $\sigma_{t+1}^{d}$ for $d=0$ and $d=1$ are known. Let $\hat{\mathbb{X}}_{t}$, $\hat{\mathbb{V}}_{t}$ and $\hat{\mathbb{X}}_{t}^{\prime}$ be sequences of points satisfying the following Euler equation for workers:
%
\begin{equation}\label{eq: Euler_func_ret}
    u^{\prime}((1+r)\hat{x}_{i} + dy - \hat{x}^{\prime}_{i}) = \beta (1+r) u^{\prime}((1+r)\hat{x}^{\prime}_{i} + y d_{t+1} - \sigma_{t+1}^{d_{t+1}}(\hat{x}^{\prime}_{i},d_{t+2})),
\end{equation}
\vspace{-1em}
\begin{equation}
    \hat{v}_{i} = u((1+r)\hat{x}_{i} + dy - \hat{x}^{\prime}_{i}) - d\tau + V_{t+1}^{d}(\hat{x}_{i}),
\end{equation}
%
where $d_{t+2} = \mathcal{I}_{t+1}(\hat{x}^{\prime}_{i},d_{t+1})$, $\hat{x}_{i}\in \hat{\mathbb{X}}_{t}$, $\hat{v}_{i}\in \hat{\mathbb{V}}_{t}$ and $\hat{x}_{i}^{\prime}\in \hat{\mathbb{X}}_{t}$. Such a sequence of points can be generated analytically using EGM. In particular, we have
%
\begin{equation}
 (1+r)\hat{x}_{i} + dy - \hat{x}^{\prime}_{i}  =  u^{\prime,-1}\left[\beta (1+r) u^{\prime}((1+r)\hat{x}^{\prime}_{i} + yd_{t+1} - \sigma_{t+1}^{d_{t+1}}(\hat{x}^{\prime}_{i},d_{t+2}))\right].
\end{equation}
%
%
In the case of the EGM, $\hat{\mathbb{X}}_{t}$ is the endogenous grid of points, $\hat{\mathbb{X}}_{t}^{\prime}$ is the exogenous one, and $\hat{\mathbb{V}}_{t}$ is the value correspondence.  

Next, order the points in $\hat{\mathbb{X}}_{t}$, $\hat{\mathbb{V}}_{t}$ and $\hat{\mathbb{X}}_{t}^{\prime}$ in ascending order of the \textit{endogenous grid points} $\hat{\mathbb{X}}_{t}$. Consider the left panel of Figure \ref{fig: FUES_1} as a stylised plot of the endogenous grid points $\hat{\mathbb{X}}_{t}$ and associated continuation payoffs $\hat{\mathbb{V}}_{t}$ generated by EGM. Assume for the purpose of illustration that the EGM points are associated with two overlapping future choice-specific value functions. The right panel of Figure \ref{fig: FUES_1} displays the policy functions (i.e., next period assets) associated with each future choice-specific value function. Pick a point $\hat{x}_{i}$, with $\hat{x}_{i}\in \hat{\mathbb{X}}_{t}$ such that $\hat{x}_{i}^{\prime}$ is optimal given $\hat{x}_{i}$ and it lies on the upper envelope.  Note if the points $\hat{x}_{i+1}$ and  $\hat{x}_{i+1}^{\prime}$ imply a different future sequence of discrete choices to $\hat{x}_{i}$ and $\hat{x}_{i}^{\prime}$, then $\hat{x}_{i+1}^{\prime}$ will experience a `discontinuous jump' from $\hat{x}_{i}^{\prime}$. However, for $\hat{x}_{i+1}^{\prime}$ to be on the upper envelope, it must be that $\hat{x}_{i+1}^{\prime}$ can only jump if it occurs after the crossing point between two value functions (for instance, the point $\hat{x}_{6}$). That is, $(\hat{x}_{i+1}^{\prime}, \hat{v}_{i+1})$ can only jump if it makes a \textit{convex} `left turn' from the line joining $(\hat{x}_{i-1}^{\prime}, \hat{v}_{i-1})$ and $(\hat{x}_{i}^{\prime}, \hat{v}_{i})$. On the other hand, if $(\hat{x}_{i+1}^{\prime}, \hat{v}_{i+1})$ makes a \textit{concave} `right turn', it cannot jump for it to be on the upper envelope. The reason is that if $(\hat{x}_{i+1}^{\prime}, \hat{v}_{i+1})$ has made a right turn, for $(\hat{x}_{i+1}^{\prime}, \hat{v}_{i+1})$ to be on the upper envelope, it must be on the \textit{concave value function} yielding the same future sequence of discrete choices $\textbf{d}$ as implied by $(\hat{x}_{i}^{\prime}, \hat{v}_{i})$. If a right turn is associated with a jump (e.g., point $\hat{x}_{7}$), then it must be on a value function associated with a sub-optimal set of future discrete choices. Section \ref{sec: theory} formally proves this argument (and provide the pseudo-code), while informally we use the intuition from Figure \ref{fig: FUES_1} to implement FUES as:
%
\begin{tcolorbox}[
colframe=black!100,
colback=white!10,
coltitle=white!100!black,  
fonttitle=\bfseries,
adjusted title= Box 1: FUES method]
\begin{enumerate}
    \item Compute $\hat{\mathbb{X}}_{t}$, $\hat{\mathbb{V}}_{t}$ and $\hat{\mathbb{X}}_{t}^{\prime}$ using standard EGM.
    \item Set a pre-determined `jump detection' threshold $\bar{M}$.
    \item Sort all sequences in order of the \textit{endogenous} grid $\hat{\mathbb{X}}_{t}$.
    \item Start from point $i= 2$. Compute $g_{i} = \frac{\hat{v}_{i} - \hat{v}_{i-1}}{\hat{x}_{i}- \hat{x}_{i-1}}$ and $g_{i+1} = \frac{\hat{v}_{i+1} - \hat{v}_{i}}{\hat{x}_{i+1}- \hat{x}_{i}}$.
    \item If $\vert \frac{\hat{x}^{\prime}_{i+1} - {x}^{\prime}_{i}}{\hat{x}_{i+1} - \hat{x}_{i}}\vert >\bar{M}$ and a right turn is made ($g_{i+1}< g_{i}$), then remove point $i+1$ from grids $\hat{\mathbb{X}}_{t}$, $\hat{\mathbb{V}}_{t}$ and $\hat{\mathbb{X}}_{t}^{\prime}$. Otherwise, set $i= i+1$.
    \item If $i+1\leq \vert \hat{\mathbb{X}}_{t}\vert$, then repeat from step 5. 
\end{enumerate}
\end{tcolorbox}
%
Note that the `jump detection' threshold $\bar{M}$ arises naturally in economic problems since the maximum curvature of a policy function \textit{conditional on a  sequence of future discrete choices} is bounded above as marginal propensities to save are also bounded above \citep{White2015, Carroll2023}. Thus, $\bar{M}$ can simply be set to the maximum marginal propensity to save, or it can also be analytically evaluated for each grid point (see Appendix \ref{sec:ENDJTAppx}). 

%tikzfigures 2 and 3
\begin{figure}[H]
\begin{minipage}[t]{0.45\textwidth}
\centering
\scalebox{0.45}{\input{figures/tikzfig2.tex}}
\end{minipage}
\hspace{0.50cm}
\begin{minipage}[t]{0.45\textwidth}
\centering
\scalebox{0.45}{\input{figures/tikzfig3.tex}}
\end{minipage}
% Common X-axis label
\begin{center}
\smallskip
\hspace{-1.5cm}
\text{\small Endogenous grid of financial assets at time t}
\end{center}
% Common X-axis label (manually placed)

\caption{FUES eliminates points that cause a concave `right turn' (left panel) from an optimal point and also cause a discontinuous jump in policy (right panel).}
\label{fig: FUES_1}
\end{figure}
%
The method will thus yield a set of refined grids $\mathbb{X}_{t}$, $\mathbb{V}_{t}$ and $\mathbb{X}_{t}^{\prime}$. To address the occasionally binding financial assets constraint $a_{t+1}\geq 0$, we follow the inverse Euler equation approach formalized by \cite{DS2024RFC}. In particular, we construct a constrained region of solution points as an array of consumption values $c_{t}^{1}$ such that $u^{\prime}(c_{t}^{1})> \beta(1+r)u^{\prime}(c^{\prime})$, where $c^{\prime}$ is next period consumption if $a^{\prime}=0$. The constrained region can then be appended to the unconstrained endogenous grid to create a single endogenous grid. The policy and value functions can then be interpolated over these grids to yield the time $t$ approximated solution for the worker who chooses to continue working. The retiree value and policy functions can be calculated using standard EGM since the retiree problem is concave at each time $t$. Once the retiree problem is solved for time $t$, the discrete choice at time $t$ of whether or not to work at time $t+1$ can be evaluated. The procedure can then be repeated at $t-1$ as per standard backward policy iteration. 

%%%%%%%%%%%%%%%%%%%

%\subsubsection{Forward and backward scans and crossing points}\label{sec: fwd_scan}

\paragraph{Forward and backward scans and crossing points}

We now detail an additional step of FUES that allows for accuracy around kink points.\footnote{For higher grid sizes, we find that the forward and backward scans do not significantly improve accuracy. Moreover, since these additional steps add at best a first order approximation of the crossing points, the approximation error they will imply will not be an improvement on the standard linear approximation error. However, for consistency with our main theoretical results, the forward and backward scans are  part of our standard implementation of FUES.} So far, we assumed that if a point $(\hat{x}_{i+1}, \hat{v}_{i+1})$ makes a left turn from point $(\hat{x}_{i}, \hat{v}_{i})$, then it is a sufficient and \textit{necessary} condition for the point $(\hat{x}_{i+1}, \hat{v}_{i+1})$ to lie after a crossing point of two value functions between itself and $(\hat{x}_{i}, \hat{v}_{i})$. However, a point $(\hat{x}_{i+1}, \hat{v}_{i+1})$ may lie after such a crossing point and yet not generate a left turn with respect to point $x_{i}$ - see the left panel of Figure \ref{fig: fscan}. We also assumed that the first point after a crossing point $(\hat{x}_{i+1}, \hat{v}_{i+1})$ must be on the optimal choice-specific value function. However, a selected point may be the first point after a crossing between choice-specific value functions and yet not be optimal - see the right panel of Figure \ref{fig: fscan} where the point $(\hat{x}_{i}, \hat{v}_{i})$ is sub-optimal but will not be removed by the basic FUES method since it is on the same sequence of future discrete choices as the optimal point right before the crossing.
%
\begin{figure}
\centering
\scalebox{0.95}{\input{figures/scan_2.tikz}}
\captionsetup{width=.9\linewidth}
\caption{Forward and backward scans to improve accuracy of FUES around crossing points.}\label{fig: fscan}
\end{figure}
%
To rectify the issue in the left panel of Figure \ref{fig: fscan}, we can implement a forward scan before a point is eliminated. A forward scan picks a point $\hat{x}_{q}$ to the right of $\hat{x}_{i+1}$ that is on the same value function as $\hat{x}_{i}$. The point $\hat{x}_{q}$ can be found by checking points after $i+1$ in the endogenous grid and finding the first point whose policy function does not jump from the point $\hat{x}_{i}$.  We can then check to see if point $\hat{x}_{i+1}$ dominates the segment joining $(\hat{x}_{i},\hat{v}_{i})$ to $(\hat{x}_{q}, \hat{v}_{q})$, drawn as $g_{w}$ in the left panel of Figure \ref{fig: fscan}. If $\hat{x}_{i+1}$ dominates the segment, it means $\hat{x}_{i+1}$ lies after a crossing point and must be included as an optimal point.

To rectify the issue in the right panel of Figure \ref{fig: fscan}, we can implement a backward scan after making a left turn. A backward scan picks a point $\hat{x}_{q}$ to the left of $\hat{x}_{i}$ that is on the same value function as $\hat{x}_{i+1}$. The point $\hat{x}_{q}$ can be found by checking points before $i$ in the endogenous grid and finding the first point whose policy function does not jump from the point $\hat{x}_{i+1}$.  We can then check to see if point $\hat{x}_{i}$ is dominated by the segment joining $(\hat{x}_{i+1},\hat{v}_{i+1})$ to $(\hat{x}_{q}, \hat{v}_{q})$, drawn as $g_{w}$ in the right panel of Figure \ref{fig: fscan}. If $\hat{x}_{i}$ is dominated by the segment, it means $\hat{x}_{i}$ lies after a crossing and must not be included as an optimal point.

Finally, note that the forward and backward scan procedures can be used to attach approximations of the choice-specific value function crossing points to the endogenous grid. For instance, in the right panel of Figure \ref{fig: fscan}, a crossing point can be attached at the intersection of the segment $g_{w}$ with the segment $g_{2}$.


%%%%%%%%%%%%%%%%%%%

\subsubsection{Challenges with existing methods}\label{subsubsec:challenges}

Before turning to the numerical results, we briefly review the main 
upper envelope algorithms in the literature, outline their practical limitations, and
explain how FUES overcomes these drawbacks. For completeness, we provide detailed pseudo-code and further discussion of \citet{Iskhakov2017} and
\citet{Druedahl2017} in Appendix~\ref{appx:existing}.\footnote{Our focus remains on upper-envelope methods, and thus on \citet{Iskhakov2017} and
\citet{Druedahl2017}. To our best knowledge, \citet{Fella2014} was the first to detail the
multiple optima problem in discrete-continuous EGM. However, the proposed solution relies on local
numerical optimization in regions where choice-specific value functions
overlap, rather than on a dedicated upper envelope envelope routine.} A key drawback of existing upper envelope methods is the inability to handle non-monotone policies. We point out that while the above illustrative example does not feature non-monotone policy functions, continuation state policies that are not strictly monotone will arise when we turn to a problem with both financial savings and continuous housing investment with fixed adjustment costs (Section \ref{sec: app2_cont_hous}) and a problem with discrete choices and present bias (Section  \ref{sec: Fella}).\footnote{Note these problems are not within the scope the general multi-dimensional version of LTM since the state-space is still one-dimensional.} Accordingly, we provide a brief review of the non-monotonicity problem below and return to it in more detail in Sections \ref{sec: egm_app2} and \ref{sec: egm_app_3}.

\paragraph{Monotone segment selection} A prominent approach to constructing upper envelopes generated by DC-EGM problems is monotone segment selection (MSS) provided by \citet{Iskhakov2017}. MSS exploits the fact that if the continuation state policy function is strictly monotone, then monotone overlapping segments that lead to multiple optima can be separately identified by detecting where the endogenous grid decreases along the exogenous grid. Within each such segment, a single future sequence of discrete choices is implied. The algorithm interpolates
the value function on every segment separately and then evaluates all interpolants on a common grid to assemble the upper envelope. 

While \cite{Iskhakov2017} demonstrate via a numerical application to the retirement choice problem that MSS can successfully recover the upper envelope,  its first drawback is that its complexity increases as the number of overlapping future choice-specific value functions grows. This is because as the number of jumps grows, there will be a larger number of distinct segments that must be identified, stored as a separate objects and then stitched together. Moreover, the simple retirement choice problem we presented above features monotone continuation state policies (the policy function can only jump `up'). If the continuation state policy is non-monotone policies, then EGM will furnish multiple endogenous grid point in the state-space for a given exogenous grid point in the continuation state-space. In this case, distinct overlapping segments cannot be identified (See \cite{Iskhakov2017}, p339) since endogenous grid points belonging to distinct segments may increase or decrease from preceding points along the exogenous grid.

\paragraph{Local triangulation method} A second approach effectively computes the upper envelope by identifying and comparing local interpolants generated by overlapping future choice specific value functions \citep{Druedahl2017,Druedahl2021}.  Motivated by its
multi-dimensional extension, we refer to this strategy as the \textit{local
triangulation method} (LTM).  LTM searches the exogenous grid for intervals (pairs of ordered points) whose corresponding endogenous grid points cover a given point in the state-space grid and then takes the maximum of the corresponding interpolants. 

Like MSS, LTM requires monotone continuation state policies. This is because for LTM to succeed, intervals in the exogenous grid must be composed of consecutive points that belong to the same future discrete choice sequence (that is, there is no jump in-between the points in the exogenous grid), \textcolor{orange}{however, if an exogenous grid point admits multiple endogenous grid pints, there is no way to determine well-defined intervals.}  The LTM also sees an increase in complexity as grid size increases. In particular, the complexity of LTM increases quadratically with grid size, since each loop is performed over both the exogenous grid and state-space grid -- we turn to the complexity of LTM in more detail as we discuss the numerical results below in Section \ref{subsubsec:application1numerical}. 

\paragraph{Roof-top cut} Another recent upper-envelope approach is the multi-dimensional roof-top cut (RFC) algorithm suggested in \cite{DS2024RFC}. RFC is also a general multi-dimensional algorithm and uses the gradient of the future choice specific value function at each point and removes points below the tangent plane constructed by the gradient. While RFC is able to handle non-monotone policies, the method still relies on exhaustive nearby nearest neighbor searches, which, in the one dimensional case, may add complexity. \footnote{\textcolor{orange}{RFC uses the analytical gradient of the value function at each EGM point along with a nearest neighbor search to evaluate whether neighboring points are on the upper envelope. Here we compare FUES with RFC in its general form. While the generalized RFC may be computationally less efficient than FUES, it can be adapted for specific applications to reduce the complexity of its neighborhood search. FUES, however, maintains the advantage of requiring no such ad hoc adjustments, offering a consistently robust and user-friendly approach to solve one-dimensional models.}}

\paragraph{FUES advantage} FUES resolves both problems discussed above. First, it discards suboptimal points in a single
linear pass: a \emph{local secant--slope} test applied to the current envelope
removes dominated candidates, eliminating the need for collecting and interpolating over multiple segments or
covering interval searches.  Second, FUES directly tests whether an individual grid point produced by EGM can belong to an upper envelope, rather than construct segments or interpolate over from the exogenous grid. This imposes no restrictions on the
structure of exogenous points and admits situations where the exogenous points not form a one-dimensional grid due to non-monotonicity. In
addition, unlike any of the methods above, theoretical results in Appendix~A guarantee that FUES recovers the
true upper envelope under general conditions. 

%%%%%%%%%%%%%%%%%%%

\subsubsection{Numerical implementation}\label{subsubsec:application1numerical}

\textcolor{blue}{[AS to ID: Note all exact numbers will be automatically populated with final results once they are finalized.]}

We apply our method to solve the model studied in \cite{Iskhakov2017}, with Figure \ref{fig: ret_plot_1} showing how FUES removes sub-optimal points in the value function (left panel) and selects the optimal policy function (right panel). We set the jump detection threshold $\bar{M}$ to $1.01$ -- reflecting the maximum MPS in this model.\footnote{Experiments using the analytical jump detection threshold (see Appendix \ref{sec:ENDJTAppx}) did not yield notable accuracy gains, see the {\textcolor{blue}{online repository}}  for details.} (Figure \ref{fig: finite_plot_pol} in the Appendix plots the policy functions for workers at different ages, as a direct comparison to Figure 3 in \cite{Iskhakov2017}.) To compare numerical performance, we also compute the upper envelope using (i) RFC as in \cite{DS2024RFC}, (ii) monotone segment selection (MSS) as in \cite{Iskhakov2017} using the code libraries provided by \cite{HARK208}, and (iii) a one-dimensional version of the local triangulation method (LTM) using the \texttt{ConsumptionSaving} package \citep{druedahl2024consav}. Table \ref{tab:ret_table_uetime} shows each method's upper envelope speed and mean \textcolor{orange}{L2 norm} error (comparing the true analytical consumption function with the numerically approximated one) for different sizes of financial assets grid and utility costs $\tau$.\footnote{All methods are tested using a single core cpu and \texttt{Python 3.12}. Both FUES and LTM are \textcolor{orange}{just in time} (jit) compiled using \texttt{numba}, while the segment search for MSS is \textcolor{orange}{highly optimized using \texttt{xxx}}. All upper envelope methods are applied to a common lifecycle problem using a standardized \textcolor{orange}{upper envelope comparison class}(\texttt{UECompare}). Interactive notebooks that implement all methods and replicate the results below are available in the \href{http://github.com/akshayshanker/FUES_EGM}{\textcolor{blue}{online repository}} for this paper.}

\begin{figure}[H]
\centering
\captionsetup{width=.7\linewidth}
\hspace{-1.5cm}\includegraphics[scale = 0.75]{plots/ret_vf_aprime_all_17_3000.png}
  \captionof{figure}{Value correspondence and optimal points for $t=17$.}
  \label{fig: ret_plot_1}
\end{figure}

All four methods in Table \ref{tab:ret_table_uetime} attain essentially identical accuracy: the \textcolor{orange}{L2 norm errors hover around $-1.54$ for the baseline} calibrations and drifts only slightly at the most irregular grids.
The real distinction in this application is speed. FUES is always fastest, followed by MSS, then RFC, with LTM a distant fourth.

Across grid sizes, FUES is between \textcolor{orange}{one and three orders of magnitude faster} than LTM and between \textcolor{orange}{one to fourteen times faster} than the best alternative EGM-based scheme (MSS).
For example, with $1,000$ grid points and $\tau=1$ the envelope step completes in 0.36 milliseconds (ms) under FUES, versus 2.66ms for MSS, 2.87ms for RFC, and 9.81ms for LTM.
Refining the grid to $3,000$ points increases the FUES time only to 0.47ms, whereas MSS, RFC and LTM rise to 6.82ms, 11.15ms and 83.81s, respectively.
Thus FUES is about 14 times faster than MSS and roughly 180 times faster than LTM on this grid.
The gap widens further on the $10,000$-point grid, where FUES still finishes in under one millisecond while LTM requires more than 900ms -- a speedup factor of \textcolor{orange}{over one thousand.} 

These patterns confirm the predicted complexity differences. Let $N^{\mathrm a}$ be the number of EGM asset nodes and $N^{\mathrm m}$ the number of cash-on-hand nodes on the common evaluation grid.
The double loop in the LTM upper-envelope routine executes $O(N^{\mathrm a} N^{\mathrm m})$ operations; when both grids are refined together ($N^{\mathrm a}\propto N^{\mathrm m}\propto K$) the workload is $O(K^{2})$.
FUES avoids the inner loop altogether, so its runtime grows sub-linearly, while MSS grows roughly linearly and LTM remains quadratic.\footnote{See the upper envelope routine \texttt{upperenvelope} in the \texttt{ConsumptionSaving} package \citep{druedahl2024consav}; the outer loop runs over asset pairs and the inner loop over all evaluation-grid points.} 

FUES is also robust to changes in parameter values.
Its timing is virtually flat in $\tau$, whereas MSS slows markedly when larger $\tau$. In Appendix \ref{appx:retShocks}, we also test FUES in a version of the model above where the discrete choices are affected by taste shocks. While FUES recovers the upper envelope and its timing and performance remains robust, the stochastic shocks create many short decreasing segments that increase the complexity of MSS by forcing MSS to evaluate extra interpolants. In turn, this increases the speed advantage of FUES over MSS by a factor of two on average.

\input{tables/retirement_1}

Finally, to gauge the impact on overall solution cost, Table \ref{tab:total_timing_comparison} reports the time to solve the full retirement model (worker plus retiree).
Because EGM itself is efficient, the envelope step dominates total runtime, so the same ranking emerges.
With 1,000 grid points and $\tau=1$ the entire model solves in 0.42ms under FUES, compared with 2.72ms for MSS, 2.94ms for RFC and 9.88ms for LTM.
On the $3,000$-point grid FUES still finishes in 0.61s, while MSS, RFC and LTM require 6.98ms, 11.32ms and 83.96ms.
Hence FUES reduces total computation time by roughly one order of magnitude relative to the fastest alternative EGM method and by upto three orders of magnitude relative to LTM on fine grids.

\input{tables/retirement_3}

%%%%%%%%%%%%%%%%%%%%%%%%%%%%%%%%%%%%%%%%%%

\subsection{Application 2: Continuous housing investment with frictions}\label{sec: app2_cont_hous}

We now turn to demonstrate to a housing frictions model where the MSS method of \cite{Iskhakov2017} and the LTM method of \cite{Druedahl2017} fails due to the non-monotonicity of the policy function. Interestingly, this model features two dimensions of savings, one for liquid (financial) assets and one for illiquid (housing) assets. Despite being a veritable workhorse in the housing frictions literature \citep{Kaplan2014, Yogo2016a, Dobrescu2022}, the conditions that \cite{Iskhakov2015} requires for multidimensional `pure EGM' do not hold here, with the set of Euler equations not analytically invertible. This example thus also shows how one-dimensional EGM and FUES can be practically applied to a multidimensional setting where `pure EGM' fails. The solution strategy will involve numerically obtaining all the roots of the financial assets Euler equation along a single dimension, and then evaluate the optimal non-monotonic housing policy, and implicitly the optimal financial assets policy, using EGM and FUES.

% Check the clain in the footnote above. How can non-monotonicity arise in a time separable case? 

%%%%%%%%%%%%%%%%%%%

\subsubsection{Model environment}

Let non-negative financial assets be denoted by $a_{t}$, and non-negative housing assets be denoted by $H_{t}$. Consider financial assets earn a rate of return $r$, and assume for simplicity that housing earns no returns. Investments can be made in and out of the stock of $a_{t}$ without friction, while adjusting housing to a value $H_{t+1}$ requires a payment of $\tau H_{t+1}$, with $\tau>0$. In each period $t$, with $t= 0, 1\dots, T$, the agent consumes non-housing goods $c_{t}$, and invests a total of $H_{t+1}$ and $a_{t+1}$ in housing and financial assets, respectively. The agent also makes a discrete choice $d_{t}$, where investment in and out of housing assets can only be made if $d_{t} = 1$; otherwise if $d_{t}=0$ then $H_{t+1} = H_{t}$. Finally, in each period, the agent earns a stochastic wage $y_{t}$ and we assume  $(y_{t})_{t=0}^{T}$ is a finite, Markov process. 

The following constraints will hold for each $t$: first, total investments and consumption cannot exceed total available wealth each period,
%
\begin{equation}\label{eq: fric_BC1}
    (1+r)a_{t} + y_{t} +  d_{t}H_{t}  \geq a_{t+1} + c_{t} + d_{t}(1+\tau)H_{t+1},
\end{equation}
%
where total wealth at time $t$ is denoted by $x_{t} =  (1+r)a_{t} + y_{t} +  d_{t}H_{t}$.

Second, in terms of payoffs, the agent lives up to time $T$, after which they die and value the bequest they leave behind according to a function $\theta \colon \mathbb{R}_{+} \rightarrow \mathbb{R}\cup\{-\infty\}$.  Per-period utility is given by a real-valued function $u \colon \mathbb{R}_{+}\times \mathbb{R}_{+}\rightarrow \mathbb{R}\cup\{-\infty\}$ as $u(c_{t}, H_{t+1})$, where $u$ is a concave, jointly differentiable, increasing function of non-housing consumption $c_{t}$ and end-of-period housing $H_{t+1}$ \citep{Yogo2016a}. Formally, the agent's dynamic optimization problem becomes
%A
\begin{equation}\label{eq: seq_val}
   V_{0}(a_{0}, H_{0}, y_{0}) = \quad \max_{(a_{t},H_{t}, d_{t})_{t=0}^{T+1}}\left\{ \sum_{t=0}^{T} \beta^{t}\mathbb{E}u(c_{t}, H_{t+1}) + \mathbb{E}\beta^{T+1}\theta((1+r)a_{T+1})\right\},
\end{equation}
%
% AS TO AS: SHOULD THERE BE AN E ABOVE BEFORE THETA. move the expectation outside?
%
such that \eqref{eq: fric_BC1} holds, $H_{t}\geq 0$ and $a_{t}\geq 0$ for each $t$, $a_{0}, H_{0}$ and $y_{0}$ are given, and the expectation is taken over the wage process $(y_{t})_{t=0}^{T}$. 

The period $T+1$ value function is given by the bequest value $V_{T+1}(a_{T+1}, H_{T+1}) = \theta((1+r)a_{T+1} +  H_{T+1})$. For each $t\leq T$, the sequential problem implies the following recursive Bellman equation:
%
\begin{equation}\label{eq: Bellman_app2}
    V_{t}(a,H,y) = \max_{a^{\prime},H^{\prime}, d}\,\,\left\{ u(c,H^{\prime}) + \beta\mathbb{E}_{y}V_{t+1}(a^{\prime},H^{\prime},y^{\prime})\right\},
\end{equation}
%AS TO AS should we put a y in the T+1 value function?
where the prime notation indicates next period state values satisfying the budget constraint \eqref{eq: fric_BC1}, $c =  y_{t} + (1+r)a + dH - a^{\prime} - d(1+\tau)H^{\prime}$, and expectations are now conditional on the realization $y$ of the time $t$ wage $y_{t}$. 

%One final remark. In contrast to Application 1, the application here features uncertainty. The solutions to the sequential problem in Equation \eqref{eq: seq_val} provide the paths of the discrete choices. Within a path, each period's discrete choice depends on the history of shocks up to that time. However, we assume the exogenous shock takes on finitely many values each period. There are thus finitely many possible future stochastic sequences of history-dependent discrete choices that can be collected in a finite set $\DChSet$, similar to Application 1 (see discussion of Equation \eqref{eq: Bell_Q}). As a result, the intuition related to what causes the secondary kinks we discussed via Figure \ref{fig: concav_env} carries over to the context of this application too.

%%%%%%%%%%%%%%%%%%%

\subsubsection{The Euler equations}

The problem will feature two Euler equations, one for each state. For periods prior to the terminal one, the Euler equation for the financial assets is
%
\begin{equation}\label{eq: suff_cond_liq1}
    u_{1}(c_{t},  H_{t+1})\geq \beta (1+r)\mathbb{E}_{t}u_{1}(c_{t+1}, H_{t+2}),
\end{equation}
%
where we have used the subscript `$1$' to refer to the first partial derivative of $u$, and will use the subscript `$2$' to refer to the second partial derivative of $u$. If $d_{t}=1$, the Euler equation for housing is
%
\begin{equation}\label{eq: suff_cond_illiq1}
    (1+\tau)u_{1}(c_{t},  H_{t+1})\geq \overbrace{\mathbb{E}_{t}\sum_{k = t}^{\iota-1}\beta^{t-\iota}u_{2}(c_{k},  H_{k+1})}^{\substack{\text{Marginal value of} \\ \text{housing services stream}}} + \underbrace{\mathbb{E}_{t}\beta^{t-\iota}\left(u_{1}(c_{\iota},  H_{\iota+1})\right)}_{\substack{\text{Marginal value of} \\ \text{liquidating housing at time $\iota$}}}.
\end{equation}
%
The intuition of the Euler equation \eqref{eq: suff_cond_liq1} for financial assets is standard. The Euler equation for housing \eqref{eq: suff_cond_illiq1}, however, features a stochastic time subscript $\iota$, defined as the next period when $d_{\iota} = 1$. Since the next time housing is adjusted will be stochastic, $\iota$ becomes a random stopping time. The Euler equation for housing then tells us that the shadow value (price) of investment (or withdrawal) from the stock of housing assets is given by the discounted expected value of housing \emph{when housing is next liquidated}, along with the stream of housing services provided up to the time of liquidation. 
 
\paragraph{Functional Euler equations.} Since the solution sequence for the problem is recursive, there exists measurable functions $\sigma^{a}_{t}$, $\sigma^{H}_{t}$ and $\mathcal{I}_{t}$ such that $H_{t+1} = \sigma^{H}_{t}(a_{t}, H_{t}, y_{t})$, $a_{t+1} = \sigma^{a}_{t}(a_{t}, H_{t}, y_{t})$ and $d_{t}= \mathcal{I}_{t}(a_{t}, H_{t}, y_{t})$ for each t. Let $\sigma^{a,d}_{t}$ and $\sigma^{H,d}_{t}$ be the choice-specific policy functions conditional on the time $t$ discrete choice $d\in \{0,1\}$. Inserting the policy functions back into \eqref{eq: suff_cond_liq1} - \eqref{eq: suff_cond_illiq1} yields the following functional Euler equation for housing:
%
\begin{multline}\label{eq: suff_cond_liq2}
    u_{1}(c, H^{\prime})(1+ \tau) \geq \beta \mathbb{E}_{y}\mathcal{I}_{t+1}(a^{\prime}, H^{\prime}, y^{\prime})u_{1}(c^{\prime}, H^{\prime\prime}) + u_{2}(c, H^{\prime}) \\ + (1-\mathcal{I}_{t+1}(a^{\prime}, H^{\prime}, y^{\prime}))\beta \left[\mathbb{E}_{y}\Theta_{t+1}(a^{\prime}, H^{\prime}, y^{\prime}) + u_{2}(c^{\prime}, H^{\prime})\right],
\end{multline}
%
where we have that
\begin{multline}\label{eq: c_def_1}
    c^{\prime} = (1+r)\sigma^{a,1}_{t}(a,H,y) + \sigma^{H,1}_{t}(a,H,y) + y^{\prime} \\ - \sigma^{a,1}_{t+1}(a^{\prime},H^{\prime},y^{\prime}) - \sigma^{H,1}_{t+1}(a^{\prime},H^{\prime},y^{\prime})(1+\tau),
\end{multline}
%
\begin{equation}\label{eq: c_def_2}
        c = (1+r)a + H + y - \sigma^{a,1}_{t}(a,H,y) - \sigma^{H,1}_{t}(a,H,y)(1+ \tau),
\end{equation}
%
\begin{equation}\label{eq: c_def_3}
        H^{\prime\prime} = \sigma^{H,1}_{t+1}(a^{\prime},H^{\prime},y),
\end{equation}
%
$a^{\prime} =\sigma^{a,1}_{t}(a,H,y)$, $H^{\prime} =\sigma^{H,1}_{t}(a,H,y)$, and $\Theta_{t+1}$ is a multiplier denoting the continuation marginal value of housing if time $t+1$ housing stock is not adjusted. For a set of time $t+1$ recursive policy functions, we can compute $\Theta_{t}$ as a function of the states as follows:
%
\begin{multline}
    \Theta_{t}(a,H,y) = \mathbb{E}_{y}\mathcal{I}_{t+1}(a^{\prime},H^{\prime}, y^{\prime})  + u_{2}(c, H^{\prime}) \\ + \beta \mathbb{E}_{y}(1- \mathcal{I}_{t+1}(a^{\prime}, H^{\prime}, y^{\prime}))\left[\Theta_{t+1}(a^{\prime}, H^{\prime}, y^{\prime}) + u_{2}(c^{\prime}, H^{\prime}) \right].
\end{multline}

%Check above the primes take into account the correct current period choice d. How does non-monotonicity arise

The functional Euler equation for the financial assets then becomes
%
\begin{equation}\label{eq: suff_cond_liq3}
    u_{1}(c, H^{\prime}) \geq \beta (1+r) \mathbb{E}_{y}u_{1}(c^{\prime}, H^{\prime\prime}),
\end{equation}
%
with $c$, $c^{\prime}$ $H^{\prime}$  defined analogously to \eqref{eq: c_def_1} and \eqref{eq: c_def_2}, and where if the time $t$ discrete choice is not to adjust, then $d=0$. 

%%%%%%%%%%%%%%%%%%%

\subsubsection{Computation using EGM and FUES}\label{sec: egm_app2}

Suppose we know $V_{t+1}$, $\{\sigma_{t+1}^{a, d},\sigma_{t+1}^{H,d}\}_{d\in \{0,1\}}$, $\mathcal{I}_{t+1}$ and $\Theta_{t+1}$. We can first apply standard EGM with FUES to evaluate $\sigma_{t}^{a,0}$ for non-adjusters since only one Euler equation (i.e., equation \eqref{eq: suff_cond_liq3}) will hold. For each possible time $t$ housing state $H_{t}$ in the housing grid and exogenous state $y_{t}$, we can approximate $\sigma_{t}^{a,0}(\cdot, H_{t}, y_{t})$ by first setting an exogenous grid of $a^{\prime}$ values (holding $H^{\prime} = H_{t}$ fixed since no adjustment takes place), and then creating an endogenous grid of time $t$ financial assets using \eqref{eq: suff_cond_liq3} and \eqref{eq: fric_BC1}. FUES can then be applied as in Section \ref{sec: FUES_RET} to eliminate sub-optimal points, allowing us to obtain an approximation of the policy function $\sigma_{t}^{a,0}$.\footnote{To address occasionally binding constraints, we follow the approach in \cite{DS2024RFC}.}

Let us now calculate the housing policy functions for adjusters. First, note we can evaluate consumption today $c$, for a given value of $H^{\prime}$ and $a^{\prime}$, as
%
\begin{equation}\label{eq: cons_liq}
c = u_{1}^{-1}((1+r)\beta \mathbb{E}_{y}u_{1}(c^{\prime}, H^{\prime\prime}), H^{\prime}),
\end{equation}
%
where $ u_{1}^{-1}$ is the analytical inverse of $u_{1}$ in its first argument. We evaluate the policy functions for adjusters is as follows:
%
%
\begin{tcolorbox}[
colframe=black!100,
colback=white!10,
coltitle=white!100!black,  
fonttitle=\bfseries,
adjusted title= Box 2: FUES and EGM for adjuster policy function]
\begin{enumerate}
    \item Fix $y$, fix an exogenous grid $\mathbb{\tilde{H}}^{\prime}$ over $H^{\prime}$, and initialise empty arrays for the housing policy ($\mathbb{\hat{H}}^{\prime}$), value correspondence ($\mathbb{\hat{V}}$), financial assets policy function ($\mathbb{\hat{A}}^{\prime}$), and endogenous wealth grid ($\mathbb{\hat{X}}$). 
    \item For each $\hat{h}_{i}\in \mathbb{\tilde{H}}^{\prime}$:
    %
    \begin{enumerate}[label=(\roman*)]
    \item Evaluate the $P$ multiple roots to \eqref{eq: suff_cond_liq2} in terms of $a^{\prime}$, with $c$ evaluated by \eqref{eq: cons_liq} and $H^{\prime}$ fixed as $\hat{h}_{i}$; collect the roots in a tuple $(\hat{a}^{\prime,0}, \dots, \hat{a}^{\prime,p},  \dots \hat{a}^{\prime,P})$. 
    \item For each root in $(\hat{a}^{\prime,0}, \dots, \hat{a}^{\prime,p}, \dots,  \hat{a}^{\prime,P})$, evaluate the endogenous grid points of wealth $(\hat{x}^{0}, \dots, \hat{x}^{p}, \dots, \hat{x}^{P})$ using the budget constraint\footnote{To save computation time, the endogenous grid for adjusters is defined in terms of total wealth $x$, where $x = (1+r)a + y + H$ since once housing is liquidated, only total wealth will affect the agent's decision.}
    %
    \begin{equation*}
    \hat{x}^{p} = \hat{a}^{\prime, p} + (1+\tau) \hat{h}_{i} + c,
    \end{equation*}
    %
    and evaluate the current period values $(\hat{v}^{\prime,0}, \dots, \hat{v}^{\prime,p}, \dots  \hat{v}^{\prime,P})$ as
    \begin{equation*}
    \hat{v}^{\prime,p} = u(c, \hat{h}_{i}) + \mathbb{E}_{y}V_{t+1}(\hat{a}^{\prime,p}, \hat{h}_{i}).
    \end{equation*}
    %
    \item Append the $P$ multiple roots $(\hat{a}^{\prime,0}, \dots, \hat{a}^{\prime,p}, \dots \hat{a}^{\prime,P})$ to $\mathbb{\hat{A}}^{\prime}$, the time $t$ values to $\mathbb{\hat{V}}$, the endogenous grid points to $\mathbb{\hat{X}}$, and $P$ copies of $\hat{h}_{i}$ to $\mathbb{\hat{H}}^{\prime}$. 
    \end{enumerate}
    %
    \item Apply FUES (Box 1, Section \ref{sec: FUES_RET}) to the grids  $\mathbb{\hat{V}}$, $\mathbb{\hat{X}}$  and $\mathbb{\hat{H}}$ to recover the refined endogenous grid points $\mathbb{X}$, value grid $\mathbb{V}$, and policy grids $\mathbb{A}^{\prime}$ and $\mathbb{H}^{\prime}$.
\end{enumerate}
\end{tcolorbox}
                
We can apply the above steps to each $y$ in the exogenous shock grid.  With $\sigma_{t}^{H,1}$ and $\sigma_{t}^{a,1}$ approximated on a uniform grid, we can then construct $\sigma_{t}^{H}$ and $\sigma_{t}^{a}$ in the standard way, by comparing the value function for adjusters and non-adjusters at each point on the uniform grid of current period states. 



\paragraph{Non-monotonicity of policy} Before we turn to the numerical example, let us discuss how the non-monotonicity of the endogenous grid points arises for the housing-adjusting policy functions. The EGM above generates an exogenous grid $\hat{\mathbb{H}}^{\prime}$ and an endogenous grid $\hat{\mathbb{X}}$ that satisfy
%
\begin{multline}\label{eq: EGM_FOC1}
u_{1}(\hat{c},\hat{H}^{\prime}) =  \beta \mathbb{E}_{y}\mathcal{I}_{t+1}( \hat{\psi}_{t}(\hat{x},\hat{H}^{\prime}) , \hat{H}^{\prime}, y^{\prime})u_{1}(c^{\prime}, H^{\prime\prime}) + u_{2}(c, H^{\prime}) \\ + (1-\mathcal{I}_{t+1}(\hat{\psi}_{t}(\hat{x},\hat{H}^{\prime}), \hat{H}^{\prime}, y^{\prime}))\beta\mathbb{E}_{y}\left[\Theta_{t+1}( \hat{\psi}_{t}(\hat{x},\hat{H}^{\prime}) , \hat{H}^{\prime}, y^{\prime}) + u_{2}(c^{\prime}, H^{\prime})\right],
\end{multline}
%
where
\begin{equation}
\hat{c} = \hat{x} - \hat{\psi}_{t}(\hat{x},\hat{H}^{\prime}) - (1+\tau)\hat{H}^{\prime},
\end{equation}
%
and $c^{\prime}$ and $H^{\prime\prime}$ defined as in \eqref{eq: c_def_1} - \eqref{eq: c_def_3}. The function $\hat{\psi}_{t}$ gives the value of $a^{\prime}$ that satisfies the financial asset Euler equation \eqref{eq: suff_cond_liq3}, for a given value of $\hat{x}$ and $\hat{H}^{\prime}$. To see the source of the non-monotoncity of $H^{\prime}$ in terms of the $\hat{x}$  above, first note how holding $\hat{\psi}_{t}$ fixed, the term on the LHS of \eqref{eq: EGM_FOC1} falls as wealth $\hat{x}$ increases, resulting in a higher level of housing investment. However, allowing $\hat{\psi}_{t}$ to also adjust alongside $\hat{x}$ implies that $\hat{\psi}_{t}$ may experience a discontinuous jump up, causing a discontinuous rise in the marginal utility of consumption today and a discontinuous \textit{fall} in housing investment. Such a discontinuous jump in $\hat{\psi}_{t}$ occurs because $\hat{\psi}_{t}$ is not evaluated for a fixed future path of discrete choices. If the agent reaches in the future a certain level of wealth and upgrades their stock of housing, they will discontinuously decrease their non-housing consumption, causing a jump in the RHS of \eqref{eq: suff_cond_liq3}.\footnote{In terms of the monotone comparative static arguments given by \cite{Iskhakov2017} in Theorem 4., allowing $a^{\prime}$ to adjust implicitly as a function $H^{\prime}$ implies that the continuation value no longer enjoys the single crossing property since now the per-period payoff $u$ also exhibits a jump and is not concave.}

\textcolor{orange}{[1) 2-3 sentances: Add details on multiple root finding. 2) one sentence clarification on non-monotonicity (perhaps).]}

%%%%%%%%%%%%%%%%%%%

\subsubsection{Numerical implementation} 

To parameterize the model, we use the following separable specification for non-housing consumption and housing:\footnote{Implementation using a non-seperable utility function is demonstrated in the FUES online \href{http://github.com/akshayshanker/FUES_EGM}{\textcolor{blue}{repository}}.} \begin{equation} u(c,H) = \frac{c^{\gamma-1}-1}{\gamma - 1} + \alpha\log(H), \end{equation} 
%
and define bequests as $\theta(c) = \bar{\theta}u(c,0)$. For simplicity, we assume i.i.d. wage shocks with a random variable $\tilde{z}$ taking values in $[0.1, 1]$ with equal probability, and assume $y_{t} = \tilde{y}_{t}(\tilde{z})$ where $\tilde{y}_{t}$ is the wage function for females in \cite{Dobrescu2016}.

\begin{figure} \centering \captionsetup{width=.7\linewidth} \includegraphics[scale = 0.75]{plots/hous_vf_aprime_all_big_59.png} \captionof{figure}{EGM total wealth points plotted on a uniform exogenous grid in the housing model with frictions.} \label{fig: haus_bigegm} \end{figure}

Since the exogenous grid is irregular, with no clearly discernible monotone segments, neither MSS or LTM cannot be applied. However, both FUES and RFC can be used to recover the upper envelope of the points recovered from inverting the Euler equation on the exogenous grid. Thus, we benchmark FUES against RFC \citep{DS2024RFC} and NEGM \citep{Druedahl2021}. NEGM first approximates the non-adjuster policy function evaluated using EGM and DC-EGM, then solves the adjuster problem using VFI. \footnote{Since the non-adjuster policy function is monotone given $H^{\prime}$, DC-EGM can be applied to the nested non-adjuster problem. We implement the numerical optimization for NEGM using a combination of Nelder-Mead and Brent's method.}

The map of endogenous grid points to the exogenous grid $\mathbb{\hat{H}}^{\prime}$ for a 59-year-old is shown in Figure \ref{fig: haus_bigegm}. (Figure \ref{fig: fric_pol} in the Appendix further compares  policy functions computed via NEGM with those computed via FUES and EGM.) First, note that there are multiple endogenous grid points for each value of the exogenous grid. Second, the endogenous grid has no discernible monotone segments that can be identified by the MSS as in \cite{Iskhakov2017}. Third, the optimal endogenous grid `jumps down', implying the policy function is not monotone. 

The numerical performance of FUES is shown in Table \ref{table:housing1} for different sizes of housing and financial assets grid and transaction costs $\tau$. Since NEGM does not include an upper envelope step that can be directly compared to FUES, Table \ref{table:housing1} compares total iteration times for each solution method. Despite the root-finding operations performed by FUES, FUES is approximately 50\% faster than NEGM. Notably, FUES achieves an order of magnitude improvement in accuracy over NEGM on average. The performance of FUES also remains consistent across various grid sizes and $\tau$ values.

As an additional performance test, we also compute an infinite horizon version of the housing frictions model with a stationary wage process. We solve the agent's problem using backward induction until the value function errors converge to a predefined tolerance level. Table \ref{table:housinginf} in the Appendix shows that FUES maintains a similar order of magnitude gain in accuracy over NEGM. Interestingly, this better FUES accuracy coming from eliminating the numerical optimization step also resulted in an average of 20\% fewer iterations required to achieve convergence of the value function. 

\input{tables/housing_1}

%%%%%%%%%%%%%%%%%%%%%%%%%%%%%%%%%%%%%%%%%%%%%

\subsection{Application 3: Discrete housing choice with present-bias}\label{sec: Fella}

In our final application we revisit a finite horizon version of the discrete housing model of \citet{Fella2014}, embedding present bias -- formalised as quasi-hyperbolic discounting with parameters $(\beta,\delta)$ that apply an extra one-shot weight $\delta<1$ to every future utility flow. In addition, we allow households to rent rather than buy. Although present bias is central to the savings literature \citep{laibson21997,Laibson2020}, no housing model with quasi-hyperbolic agents has yet been solved with either EGM or Euler-equation techniques. With discrete housing choices the standard Euler condition need not bind: a present-biased household may wish to raise current consumption yet recognises that any reduction in financial saving jeopardises a discrete housing upgrade demanded by its future selves. The resulting consumption policy is step-shaped and non-monotone, causing all existing DC-EGM routines to fail. By contrast, FUES remains accurate, efficient, and robust in this environment.

%%%%%%%%%%%%%%%%%%%

\subsubsection{Model environment}

Consider a finitely lived  individual who draws a sequence of bounded, stationary Markov wage shocks $(y_{t})_{t=0}^{T}$ taking on values in a finite discrete space $Y\subset \mathbb{R}_{+}$. Each period the individual enters with housing stock $H_{t}$ and financial assets $a_{t}$, and decides whether they will be a renter ($H_{t+1}=0$) or home-owner ($H_{t+1}>0$). They then decide how much to consume of non-durable goods $c_{t}$, how much to save in financial assets $a_{t+1}$, and how much to save in housing assets $H_{t+1}$ if they are an owner and their housing services $H^{\text{rent}}_{t}$ if they are a renter. 

Similar to Application 2, housing assets serve a dual role, namely they are a form of investment good but also provide durable consumption services. However, housing can only be purchased in discrete amounts on a finite grid $\mathbb{H}\subset \mathbb{R}_{++}$, and adjusting housing each period incurs a fixed transaction cost $\tau H_{t+1}$, with $\tau \in [0,1)$. We assume no borrowing and so, $a_{t}\geq 0$ must hold. In what follows we will use $B_{t}\colon = \mathbb{R}_{+}\times \mathbb{H}$ to refer to the agent's start of period endogenous state, with $b_{t}\in B_{t}$ and $b_{t} = (a_{t}, H_{t})$. 

\paragraph{Budget constraints} Conditional on entering period $t$ as a homeowner with housing stock $H_t$, the budget constraint is
\begin{equation}\label{eq:budget_owner}
a_{t+1}+c_{t} \;=\; (1+r)a_{t}+y_{t}
                  + d^{\text{adj}}_{t}H_{t}
                  -(1+d^{\text{adj}}_{t}\tau)H_{t+1}
        \;:=\; w_{t}^{\text{own}},
\end{equation}
where $d^{\text{adj}}_{t}\colon = \mathbb{1}_{\{H_{t+1}\neq H_{t}\}}$ flags housing stock adjustment and payment of the adjustment cost $\tau$. The term $ w_{t}^{\text{own}}$ refers to the owner's total wealth. If the individual instead rents in period $t$, the budget constraint becomes
%
\begin{equation}\label{eq:budget_renter}
a_{t+1}+c_{t} \;=\; (1+r)a_{t}+y_{t}
                  + H_{t}
                  - P^{r}S_{t}
        \;:=\; w_{t}^{\text{rent}},
\end{equation}
%
where $S_{t}$ is the flow of rented housing services purchased at price $P^{r}$ and $ w_{t}^{\text{rent}}$ is the renter's total wealth.

\paragraph{Preferences and recursive problem} 
%
Housing services consumed in the period are denoted by $h_{t}$, with $h_{t}=H_{t+1}$ for owners and $h_{t}=S_{t}$ for renters, so, following \cite{Fella2014}, per-period utility is
%
\begin{equation}\label{eq:period_u}
u(c_{t},h_{t}) \;=\;
    \alpha\log c_{t} + (1-\alpha)\log\!\bigl[\kappa(h_{t}+\iota)\bigr],
    \qquad \kappa>0,\; \iota>0.
\end{equation}

Turning to the recursive time inconsistent problem solved, let $V_{t+1}$ denote the continuation value at the end of period $t$. For any start of period state $(a,y,H)$, we have
%
\begin{equation}\label{eq:continuation_value}
V_{t}(a,y,H)\;=\;
  \mathbb{E}_{t}\left\{\!\sum_{j=t}^{T}\beta^{T-j}\,
  u\!\bigl(c_{j},h_{j}\bigr) + \beta^{T-t}\theta((1+r)a_{T+1})\right\}
\end{equation}
%
where $c_{j}=\sigma^{c}_{j}(a_{j},y_{j},H_{j})$ and
$h_{j}=\sigma^{h}_{j}(a_{j},y_{j},H_{j})$ are the start of period policy functions for
consumption and housing services at time $t$ and $\theta$ is bequest function.  Proceeding backwards from the end of a period, given housing stock $H_{t+1}$, the owner solves
%
\begin{equation}\label{eq:Bellman_owner_cond}
\hat{W}_{t}^{\text{own}}(H_{t+1},y_{t},w_{t}^{\text{own}})
  \;=\;
  \max_{c_{t}}
  \Bigl\{
    u(c_{t},H_{t+1})
    +\delta\beta\,
     \mathbb{E}_{y_{t}}V_{t+1}(a_{t+1},y_{t+1},H_{t+1})
  \Bigr\},
\end{equation}
subject to $a_{t+1}+c_{t}=w_{t}^{\text{own}}$, where $\hat{W}_{t}^{\text{own}}$ is the owner's current value function. The term $\delta\in (0,1]$ captures the current period individual's present bias. Next, given a choice of rented services $S_{t}$, a renter solves
%
\begin{equation}\label{eq:Bellman_renter_cond}
\hat{W}_{t}^{\text{rent}}(w_{t}^{\text{rent}},S_{t},y_{t})
  \;=\;
  \max_{c_{t}}
  \Bigl\{
    u(c_{t},S_{t})
    +\delta\beta\,
     \mathbb{E}_{y_{t}}V_{t+1}(a_{t+1},y_{t+1},0)
  \Bigr\},
\end{equation}
%
subject to $a_{t+1}+c_{t}=w_{t}^{\text{rent}}$. The owner then chooses $H_{t+1}$ to maximize
%
\begin{equation}\label{eq:Bellman_owner}
W_{t}^{\text{own}}(a_{t},H_{t},y_{t})
  \;=\;
  \max_{H_{t+1}}
    \hat{W}_{t}^{\text{own}}(H_{t+1},y_{t},w_{t}^{\text{own}}),
\end{equation}
subject to \eqref{eq:budget_owner} and $w_{t}^{\text{own}}\ge 0$. Analogously, the renter chooses $S_{t}$ to maximize
%
\begin{equation}\label{eq:Bellman_renter}
W_{t}^{\text{rent}}(a_{t},y_{t})
  \;=\;
  \max_{S_{t}}
    \hat{W}_{t}^{\text{rent}}(S_{t},y_{t},w_{t}^{\text{rent}}),
\end{equation}
%
subject to \eqref{eq:budget_renter} and $w_{t}^{\text{rent}}\ge 0$. Finally, the decision to own or rent is obtained from
%
\begin{equation}
W_{t}(a_{t},y_{t},H_{t})
  \;=\;
  \max\bigl\{
    W_{t}^{\text{rent}}\bigl((1+r)a_{t}+H_{t},y_{t}\bigr),
    \;W_{t}^{\text{own}}(a_{t},y_{t},H_{t})
  \bigr\}.
\end{equation}
%
The structure is standard for a sophisticated present-biased agent:
current choices account for the fact that all future selves discount at
rate~$\delta\beta$, even though the agent's long-run optimizer would
prefer pure exponential discounting at rate~$\beta$.  Combined with discrete housing choices, this strategic
mismatch generates constraints that become the source of the model's non-monotone policy function, which we turn to next. 

%%%%%%%%%%%%%%%%%%%

\subsubsection{Euler equation}

Deriving the Euler equation presents a challenge because $V_{t+1}$ will not be differentiable, precluding the use of standard envelope conditions. \textcolor{orange}{However, time inconsistency implies that we also cannot directly employ variational arguments \citep{Stachurski2009,DS2024RFC}} as we can in discrete-continuous problems without present-bias. We will show below that the familiar Euler equation, which binds for strictly positive financial assets, will not even be necessary (let alone sufficient); anticipating the present-bias of future selves, the current Euler equation may hold with a strict inequality even when financial assets are strictly positive. 

To proceed, consider the time $t+1$ beginning of period states to the future sequence of discrete choices $\mathrm{f}^{\mathbf{d}}_{t+1}\colon X_{t}\rightarrow \left(\left\{0,1\right\}\times \mathbb{H}\right)^{\mathbb{Z}^{T+1-t}}$. In words, the evaluation $\mathbf{d}_{t+1} = \mathrm{f}^{\mathbf{d}}_{t+1}(x)$ produces a stochastic sequence of discrete choices implied by the solution to the problem from $t+1$ on-wards. If the individual at time $t$ chooses $x$ as their end of period state, their optimal choice of $\mathbf{d}$ will not agree with  $\mathbf{d}_{t+1}$ since they would prefer their future selves to use a discount factor of $\beta$ between periods -- however, they cannot induce their future selves to follow the optimal $\mathbf{d}$, rather they are \emph{constrained} by the choice of  $\mathbf{d}_{t+1}$ their future selves make. Formally, this means that we can write the time $t$ owner's problem as
%
\begin{equation}\label{eq:Bellman_owner_cond}
\hat{W}_{t}^{\text{own}}(H_{t+1},y_{t},w_{t}^{\text{own}})
  \;=\;
  \max_{c_{t}}
  \Bigl\{
    u(c_{t},H_{t+1})
    +\delta\beta\,
     \mathbb{E}_{y_{t}}Q^{\mathrm{f}^{\mathbf{d}}_{t+1}(x_{t+1})}_{t+1}(a_{t+1},y_{t+1},H_{t+1})
  \Bigr\},
\end{equation}
%
where $Q^{\mathbf{d}}$ is the time $t+1$ continuation value conditional on a future stochastic sequence of discrete choices. It follows that 
%
\begin{multline}\label{eq:Bellman_owner_cond}
\hat{W}_{t}^{\text{own}}(H_{t+1},y_{t},w_{t}^{\text{own}})
  \;=\;
  \max_{\mathbf{d}_{t+1}}\Biggr\{\max_{c_{t}}
  \Bigl\{
    u(c_{t},H_{t+1})
    +\delta\beta\,
     \mathbb{E}_{y_{t}}Q^{\mathbf{d}_{t+1}}_{t+1}(a_{t+1},y_{t+1},H_{t+1})
  \Bigr\},\\ s.t. \,\, \mathbf{d}_{t+1} = \mathrm{f}^{\mathbf{d}}_{t+1}(x_{t+1})\Biggr\}
\end{multline}
%
Re-expressing the mapping constraint as a state-space restriction,
%
\begin{equation}
  \mathbf d_{t+1}=f_{t+1}^{\mathbf d}(x_{t+1})
  \;\Longleftrightarrow\;
  x_{t+1}\in S(\mathbf d_{t+1}),
\end{equation}
%
%
with $S(\mathbf d_{t+1})\subset X_{t+1}$ connected (proved in Appendix A), we
can rewrite \eqref{eq:W_owner_inner} as a nested maximization
\begin{align}
  \hat W_{t}^{\text{own}}(H_{t+1},y_{t},w_{t}^{\text{own}})
  \;=\;
  \max_{\mathbf d_{t+1}}
  \Bigl\{
    \max_{\substack{c_{t}\\ x_{t+1}\in S(\mathbf d_{t+1})}}
    \bigl[
      u(c_{t},H_{t+1})
      +\delta\beta\,
       \mathbb E_{y_{t}}
         Q_{t+1}^{\mathbf d_{t+1}}(a_{t+1},y_{t+1},H_{t+1})
    \bigr]
  \Bigr\}.
  \label{eq:W_owner_nested}
\end{align}
%

Within each admissible set $S(\mathbf d_{t+1})$ the inner problem is smooth,
so the Kuhn-Tucker conditions yield\footnote{%
  Gradient terms involving $\partial a_{t+1}/\partial c_{t}$ have been
  substituted using the budget constraint.}
  %
\begin{equation}\label{eq:Euler_gen}
  u_{c}(c_{t},H_{t+1})
  \;\ge\;
  \mathbb{E}R\bigl[
    \frac{\partial\sigma^{\mathbf d_{t+1}}(x_{t+1})}{\partial a}\,
    \bigl(
      u_{c}(c_{t+1},H_{t+1})
      -\delta\Phi_{t+2}(x_{t+2})
    \bigr)
    +\delta^{-1}u_{c}(c_{t+1},H_{t+1})
  \bigr],
\end{equation}
with equality whenever $a_{t+1}$ lies in the interior of
$S(\mathbf d_{t+1})$. The term
%
\begin{equation}\label{eq:lambda_presentbias_gen}
  \Phi_{t+1}(x_{t+1})
    \;=\;
   \frac{\partial\sigma^{\mathbf d_{t+1}}(x_{t+1})}{\partial a}\,
    \bigl(
      u_{c}(c_{t+1},H_{t+1}) - 
      \delta\Phi_{t+2}(x_{t+2})
    \bigr)
    +\delta^{-1}u_{c}(c_{t+1},H_{t+1})
\end{equation}
%
%
is defined recursively and captures the marginal value of relaxing the financial
asset constraint one period ahead. The term on the RHS above is the standard marginal value of financial saving in a standard time inconsistent Euler equation (recall \textcolor{orange}{xxx}). However, The above equation does not always hold with equality, with the left hand boundaries of the intervals $S(\mathbf d_{t+1})\equiv [\underbar{s}, \bar{s}]$ acting as constraints. Intuitively, at $\underbar{s}$ the current self would like to
reduce saving, but doing so would trigger a discrete jump to a lower
value-function branch as their future self would switch the sequence of discrete housing and renting choices they make. Thus the current self saves $\underbar{s}$ until their current wealth is low enough they prefer the switch to $\mathbf{d}_{t+1}$. This will lead to a non-monotone savings function with flat regions where a unique inverse does not exist. 

\textcolor{orange}{2-3 sentances possibly explaining technical details and refs to math lit.}
%% AS TO AS: Concavity even without discrete choices?

%%%%%%%%%%%%%%%%%%%

\subsubsection{Computation using EGM and FUES}\label{sec: egm_app_3}

In the presence of time inconsistency, box \ref{box: app3algo} shows how to construct the exogenous grid of points. in addition to the points where the Euler equation holds with equality, we must also add the points where the Euler equation does not bind. The resulting exogenous grid will imply multiple values of endogenous grid points for each $a^{\prime}$ value--- reflecting the fact that the policy function possess flat sections. No existing methods can be applied to recover the upper envelope this problem since there are no clearly discernible monotone segments among the exogenous set of points that satisfy the Euler equation (ruling out MSS) and the points do not form a regular exogenous grid with well-defined intervals to loop over (ruling out LTM). however, FUES can be efficiently applied off the shelf. 

\begin{tcolorbox}[
  colframe=black!100,
  colback=white!10,
  fonttitle=\bfseries,
  adjusted title=Box 3: Exogenous-grid construction for the present-bias model]

\begin{enumerate}
  \item \textbf{Fix housing.}\;
        Select a housing level $H\in\mathbb H$ (owner: $H=H_{t}$; renter: $H=0$) and income shock $y$. Initialize empty $\mathbb{X}_{t}$, $\hat{\mathbb{X}}^{\prime}_{t}$ and $\mathbb{V}_{t}$. 

  \item \textbf{Initial exogenous grid.}\;
        Create a regular financial-asset grid
        $\hat{\mathbb A}_{t}^{\prime}=\{a_{1}^{\prime},\dots,a_{N_{a}}^{\prime}\}$
        on $[0,\bar a]$.

  \item \textbf{Invert the Euler equation (equality).}\;
        For every $a_{j}^{\prime}\in\hat{\mathbb A}_{t}^{\prime}$  
        anaytically invert the generalised Euler equation \eqref{eq:Euler_gen} \emph{with equality}  
        to obtain consumption $\hat c_{j}$ and the corresponding endogenous financial asset  
        $\hat x_{j}=w^{\mathrm{own/rent}}_{t}-\hat c_{j}$.  
        Append the triple  
        $(\hat x_{j},\hat{a}_{j}^{\prime},\hat v_{j})$ to the sets $(\mathbb{X}_{t}, \hat{\mathbb{X}}_{t}^{\prime}, \mathbb{V}_{t})$.

  \item \textbf{Insert constrained points.}\;
        Whenever $\hat x_{j+1}<\hat x_{j}$  
        (a downward step in the endogenous grid),  
        append $N^{\mathrm{cons}}$ \emph{constrained} points:\par
        \hspace{1em}-- fix $\tilde x_{k}=\hat x_{j+1}$,\; $k=1,\dots,N^{\mathrm{cons}}$;\par
        \hspace{1em}-- choose consumptions $\tilde c_{k}$ decreasing from $\hat c_{j+1}$ to $0$  
        so that the strict inequality in \eqref{eq:Euler_gen} holds;store $(\tilde x_{k},a_{j+1}^{\prime},\tilde v_{k})$ \par

  \item \textbf{Assemble provisional grids.}\;
        Concatenate the constrained points to the grid and
        sort by $\hat{\mathbb X}_{t}$ (endogenous assets). Now we denote the combined arrays by  
        $\hat{\mathbb X}_{t}$, $\hat{\mathbb X}_{t}^{\prime}$, $\hat{\mathbb V}_{t}$.

  \item \textbf{Apply FUES.}\;
        Run the FUES algorithm on  
        $\bigl(\hat{\mathbb X}_{t},\hat{\mathbb V}_{t},\hat{\mathbb X}_{t}^{\prime}\bigr)$  
        to delete dominated pairs and recover the period-$t$ upper envelope.
\end{enumerate}
\end{tcolorbox}\label{box: app3algo}

%% AS TO AS: exogenous grid vs. exogenous set of points. 

\subsubsection{Numerical implementation}

\begin{figure} \centering \captionsetup{width=.7\linewidth} \includegraphics[scale = 0.5]{plots/pb_housing/egm_grid_H0_y0_period0_FUES.png} \captionof{figure}{EGM total wealth points plotted on a uniform exogenous grid in the housing model with frictions.} \label{fig: haus_bigegm} \end{figure}

Figure \ref{fig: haus_bigegm} overlays the endogenous grid points produced by FUES-EGM on the exogenous cash-on-hand grid for the frictionless benchmark ($\beta=1$). Relative to value-function iteration (VFI) the algorithm attains the same solution in roughly $x$-fold less CPU time, a gain that mirrors the order-of-magnitude speed-ups reported for RFC, MSS and LTM in Appendix \ref{app: benchmarks}.
When quasi-hyperbolic discounting is switched on the policy surface acquires pronounced plateaux: as shown in Figure \ref{fig: savings}, consumption is flat over wide wealth intervals because the household hoards financial assets purely to satisfy the down-payment constraint of its next discrete housing upgrade. The resulting stepped (and globally non-monotone) schedule breaks all existing DC-EGM routines, whose local training cannot locate the upper envelope. FUES, by contrast, identifies the correct policy in a single targeted secant pass, preserving both accuracy and numerical stability while retaining its large speed advantage over VFI.

\begin{figure}[H] \centering \captionsetup{width=.7\linewidth} \includegraphics[scale = 0.5]{plots/pb_housing/savings.png} \captionof{figure}{Savings.} \label{fig: haus_bigegm} \end{figure}


\textcolor{orange}{Non-monotonicity of housing policy discussion.}
%%%%%%%%%%%%%%%%%%%%%%%%%%%%%%%%%%%%%%%%%%%%%

\section{General formulation of FUES}\label{sec: theory}

We now formally define a general discrete-continuous optimization problem, apply the EGM and state the FUES method as a formal pseudo-code, providing an illustrative proof to show how FUES recovers the optimal endogenous grid points. 

%%%%%%%%%%%%%%%%%%%

\subsection{The general discrete-continuous optimization problem}

Let the set $\DChSet$ be a finite family of sequences of future discrete choices, and let $\left\{G^{\dVecVar}\right\}_{\dVecVar\in \DChSet}$ be a family of continuously differentiable, increasing and strictly concave functions. Each function $G^{\dVecVar}$ is indexed by a future sequence of discrete choices. Assume that $G^{\dVecVar}\colon S\times Z\rightarrow \mathbb{R}\cup\{-\infty\}$ for each $\dVecVar$, with $S\subset \mathbb{R}$ and $Z\subset \mathbb{R}$, and that $S$ and $Z$ are compact. In the context of a dynamic programming problem, each element in the choice set  $\DChSet$ will be a future stochastic sequence of history-dependent choices (and not a single discrete choice at a given time). The function $G^{\dVecVar}$ will be the maximand on the RHS of the Bellman operator, holding the discrete choice fixed. Now, consider maximizing the upper envelope $U$ of the choice-specific functions wrt. a continuous choice variable $z$
%
\begin{equation}\label{eq: problem}
    V(x) = \max_{z\in Z}\quad  U(x,z),
\end{equation}
%
where
\begin{equation}
U(x,z) =  \max_{\dVecVar\in\DChSet}\quad G^{\dVecVar}(x,z),
\end{equation}
%
 such that $x$ is fixed. Next, define $\mathcal{I}(x,z) = \argmax\limits_{d} G^{\dVecVar}(x,z)$ as the optimal discrete choice holding the arguments of each $G^{d}$ fixed. Also, define the optimal policy function as

\begin{equation}
    \sigma(x) = \argmax_{z}\quad U(x,z),
\end{equation}
%
and the choice-specific policy function and choice-specific value function as
%
\begin{equation}
    \sigma^{\dVecVar}(x) = \argmax_{z}\quad  G^{\dVecVar}(x,z),\qquad Q^{\dVecVar}(x) =  \max_{z}\quad  G^{\dVecVar}(x,z).
\end{equation}
	
Finally, define a twice differentiable and invertible transition function $f\colon S\times Z \rightarrow S$ that maps the current period choice and state to next period state in a dynamic programming problem. We set $x^{\prime}$ to denote the next period state, where $x^{\prime}$ will obey $x^{\prime}= f(x,z)$. 

Note that the mapping between our discrete-continuous optimization problem as stated above and the three applications we tackle is as follows:

\emph{Application 1, Section \ref{sec: retchoice}.} Fix $t$ and consider agents making a decision to continue work in time $t+1$. Each $\dVecVar$, with $\dVecVar\in \DChSet$, corresponds to a future path of feasible discrete choices made from $t+1$ to $T$. Formally, we can write  $\DChSet = \left\{0,1\right\}^{T-t}$. The choice-specific payoff $G^{\dVecVar}(x,z)$ corresponds to $u((1+r)x + y - z) + V_{t+1}^{\dVecVar}(z)$, where $x$ denotes financial assets, $z$ is the choice of next period assets, and $V^{\dVecVar}_{t+1}$ is the value function for those who work in $t+1$ conditional on the future path of discrete choices made from $t+1$ to $T$.

\emph{Application 2, Section \ref{sec: app2_cont_hous}.} Fix $t$, recall our discussion in Section \ref{sec: egm_app2} and consider those making a decision to adjust their housing at time $t$. Each $\dVecVar$, with $\dVecVar\in \DChSet$, corresponds to a future history-dependent path of housing adjustment choices from $t+1$ to $T$. The choice-specific payoff $G^{\dVecVar}$ corresponds to
%
\begin{equation*}
G^{d}(x, z) = u(x - \hat{\psi}_{t}^{d}(x,z) - (1+\tau)z,z) + \mathbb{E}_{y_{t}}V_{t+1}^{d}(\hat{\psi}_{t}^{d}(x,z), z, y^{\prime}),
\end{equation*}
%
where $x$ is the total (financial and housing) wealth, $z$ represents the $t+1$ housing choice, and $V^{d}_{t+1}$ is the $t+1$ choice-specific value function. For a fixed discrete choice, the function $\hat{\psi}_{t}^{\dVecVar}$ maps the wealth and housing choice to the optimal value of next period financial assets defined by the Euler equation \eqref{eq: suff_cond_liq2}. 

\emph{Application 3, Section \ref{sec: Fella}.} Fix $t$ and consider those making a decision to adjust housing at time $t$. Each $\dVecVar$, with $\dVecVar\in \DChSet$, corresponds to an infinite history-dependent path of future housing choices. The choice-specific payoff $G^{\dVecVar}$ corresponds to the continuation payoff
%
\begin{equation}
G^{\dVecVar}(x, z) = u((1+r)x + y_{t} + (1-\delta)H - z - (1+\tau)H^{\prime},H^{\prime}) + \mathbb{E}_{y^{\prime}}V^{\dVecVar, \prime}(z, y^{\prime}),
\end{equation}
%
where $z$ represents the next period financial assets choice, $V^{\dVecVar, \prime}$ is the next period choice-specific value function, and $H$ and $H^{\prime}$ are the fixed beginning- and end-of-period housing assets. 

Since we have stated the above examples as reduced form dynamic programming problems, we have $x^{\prime} = f(x,z) = z$. 

%%%%%%%%%%%%%%%%%%%

\subsubsection{The EGM and formal FUES pseudo-code}

Let $\hat{\mathbb{X}}^{\prime}$ be the exogenous grid of values $x^{\prime}$, with $\vert \hat{\mathbb{X}}^{\prime}\vert = N$, where $N$ is the exogenous grid size. We assume grids are sequences, that is $\hat{\mathbb{X}}^{\prime} = \{\hat{x}_{0}^{\prime},\hat{x}_{1}^{\prime}\dots \hat{x}_{i}^{\prime}\dots\}$. Recall that a necessary FOC for an interior solution to the discrete-continuous optimization problem is that $G^{\mathcal{I}(x,z)}_{2}(x,z) = 0$, where the subscript `$2$' denotes the partial derivative with respect to the second argument of $G^{\dVecVar}$. Also recall that if a continuous choice satisfies the FOC, it will be optimal conditional on the discrete choice (although, the discrete choice may not optimal). Formally:
%
\begin{remark}
For $x\in S$, consider $\tilde{z}$ that satisfies $G_{2}^{\mathcal{I}(x,\tilde{z})}(x, \tilde{z}) = 0$. If $\dVecVar = \mathcal{I}(x,\tilde{z})$ is fixed, then $\tilde{z}$ will satisfy $\tilde{z} = \argmax_{z\in Z} G^{\dVecVar}(x,z)$.
\end{remark}

We now define the unrefined endogenous grid of values $\hat{x}$ as
%
\begin{equation}
\hat{\mathbb{X}}\colon =  \left\{\hat{x}\in S\,\big\vert\, \hat{x}^{\prime} = f(\hat{x},z),  G^{\mathcal{I}(\hat{x},z)}_{2}(\hat{x},z) = 0, \hat{x}^{\prime}\in \mathbb{X}^{\prime}, z\in Z\right\},
\end{equation}
%
and we can also calculate unrefined grids containing candidate policies and values as
%
\begin{equation}
\hat{\mathbb{Z}}\colon =  \left\{\hat{z}\in Z\,\big\vert\, \hat{x}^{\prime} = f(x,\hat{z}),  G^{\mathcal{I}(x,\hat{z})}_{2}(x,\hat{z}) = 0, \hat{x}^{\prime}\in \mathbb{X}^{\prime}, x\in S\right\},
\end{equation}
%
\begin{equation}
\hat{\mathbb{V}}\colon =G^{\mathcal{I}(\hat{\mathbb{X}},\hat{\mathbb{Z}} )}(\hat{\mathbb{X}}^{\prime}).
\end{equation}

Note the definition of the endogenous grid values above does not necessarily imply the values of $\hat{\mathbb{X}}$ can be calculated analytically. Rather, we leave open the possibility that the endogenous grid is a general grid (computed analytically or numerically) of points that satisfy the necessary FOCs of a discrete-continuous optimization problem.

The endogenous grid will be initially ordered in ascending order according to the order of $\hat{\mathbb{X}}^{\prime}$. We then order the grids $\hat{\mathbb{X}}^{\prime}$, $\hat{\mathbb{Z}}$, $\hat{\mathbb{V}}$, and $\hat{\mathbb{X}}$ by the values in $\hat{\mathbb{X}}$, and through a slight abuse of notation, continue to denote the reordered grids as $\hat{\mathbb{X}}^{\prime}$, $\hat{\mathbb{Z}}$, $\hat{\mathbb{V}}$ and $\hat{\mathbb{X}}$. For each of these grids, the FUES method will generate a subsequence of refined grids $\mathbb{X}^{\prime}$, $\mathbb{Z}$, $\mathbb{V}$ and $\mathbb{X}$ indexed by $i_{k}$, for instance, $\mathbb{X} = \left\{\hat{x}_{0},\hat{x}_{1}, \dots, \hat{x}_{i_{k}}, \dots\right\}$. In the algorithm below, we let $N = \vert \hat{\mathbb{X}}\vert$ and $\bar{M}$ (i.e., the constant denoting the threshold at which FUES registers a jump between policy functions) be chosen by the researcher. Finally, $\mathcal{C}_{l+1}$ denotes a Boolean value that is true if point $x_{l+1}$ is optimal, and $\mathcal{D}_{l}$ denotes a Boolean value that is true if  $x_{k}$ is not optimal and should be deleted.

\vspace{-0.5em}

%The forward and backward scans also return a crossing point coordinate, which is denoted by $(\bar{x}_{l}, \bar{v}_{l})$. 
%
\RestyleAlgo{ruled}
\SetKwComment{Comment}{/* }{ */}
\begin{algorithm}[hbt!]
\caption{Fast upper envelope scan}\label{alg: NCE_scan}
\KwData{$ N > 0$, $\bar{M}>0$, $\hat{\mathbb{X}}^{\prime}$, $\hat{\mathbb{Z}}$,$\hat{\mathbb{V}}$,$\hat{\mathbb{X}}$}
\KwResult{$\mathbb{X}^{\prime}$, $\mathbb{Z}$, $\mathbb{X}$ and $\mathbb{V}$}
$k \gets 1$, $l\gets 1$\;

$\hat{x}^{\prime}_{i_{q}}\gets \hat{x}^{\prime}_{q}$, $\hat{v}_{i_{q}}\gets \hat{v}_{q}$, $\hat{x}_{i_{q}}\gets \hat{x}_{q}$ \qquad $q\leq 1$  \Comment*[r]{Init. first two points.}
%
\While{$l< N$}{

  $g^{v}_{k} \gets \frac{\hat{v}_{i_k} - \hat{v}_{i_{k-1}}}{\hat{x}_{i_{k}} - \hat{x}_{i_{k-1}}}$, $g^{v}_{l+1} \gets \frac{\hat{v}_{l+1} - \hat{v}_{i_k}}{\hat{x}_{l+1} - \hat{x}_{i_k}}$, $g^{x}_{l+1} \gets \frac{\hat{x}^{\prime}_{l+1}-\hat{x}^{\prime}_{i_k}}{\hat{x}_{l+1} - \hat{x}_{i_k}}$\;
  
  \uIf{If $g^{v}_{l+1}\leq g^{v}_{k}$ and $\vert g^{x}_{l+1}\vert > \bar{M}$ \Comment*[r]{Right turn, jump}}{$\mathcal{C}_{l+1}\gets \text{False}$\;$\mathcal{C}_{l+1}\gets$  Forward scan \Comment*[r]{Optional fwd. scan (Fig. 5)}}\label{lst:line: right_jump}
  \uElseIf{$g^{v}_{l+1}> g^{v}_{k}$\label{lst:line: left_turn} \Comment*[r]{Left turn}}{
                    $\mathcal{C}_{l+1} \gets$ True \;
                    $\mathcal{D}_{l}$ $\gets$ Backward scan \Comment*[r]{Optional bwd. scan (Fig. 5)}}
   \Else{$\mathcal{C}_{l+1} \gets$ True \label{lst:line: right_no_jump}  \Comment*[r]{Right turn, no jump}}
    \uIf{$\mathcal{D}_{l}$}{$\hat{x}^{\prime}_{i_{k}}\gets \hat{x}^{\prime}_{l+1}$, $\hat{v}_{i_{k}}\gets \hat{v}_{l+1}$\; $\hat{x}_{i_{k}}\gets \hat{x}_{l+1}$, $\hat{z}_{i_{k}}\gets \hat{z}_{l+1}$  \Comment*[r]{Replace $i_{k}$ sub-opt. point.}}
    \uElseIf{$\mathcal{C}_{l+1}$}{$\hat{x}^{\prime}_{i_{k+1}}\gets \hat{x}^{\prime}_{l+1}$, $\hat{v}_{i_{k+1}}\gets \hat{v}_{l+1}$\; $\hat{x}_{i_{k+1}}\gets \hat{x}_{l+1}$, $\hat{z}_{i_{k+1}}\gets \hat{z}_{l+1}$\Comment*[r]{Add point; keep $i_{k}$ opt. point.} $k = k+1$}
    \Else{Pass}
    $l\gets l+1$\;}
\end{algorithm}
%
%
%The main condition of the FUES method, at line 5 of Algorithm \ref{alg: NCE_scan} checks to see if a discontinuity occurs at a concave `right turn' and eliminates such candidate points. As discussed in Section \ref{sec: retchoice}, the purpose of the forward scan and backward scan is to improve approximation at the crossing point between value functions. The forward scan deals with the case when there is a point after a crossing between two choice-specific value functions and the candidate point does not make a convex `left turn' from the point before the crossing point. The backward scan deals with the case where the first point after a crossing point between two value functions is not from the `new value function'. A discussion of the Assumptions is provided after the main result. 

%%%%%%%%%%%%%%%%%%%

\newpage

\subsection{Main accuracy result}

%%%%%%%%%%%%%%%%%%%

The optimal subset of points $\mathbb{X}^{\star}$ of the endogenous grid can be defined as
%
\begin{equation}
\mathbb{X}^{\star}\colon = \left\{\hat{x}_{i}\in \hat{\mathbb{X}}\,\big\vert \,\, \hat{v}_{i} = V(\hat{x}_{i})\right\}.
\end{equation}
%

In terms of required assumptions, we first assume that the distance between choice-specific policy functions is bounded below. The first assumption is satisfied under the condition future choice specific policy functions do not intersect and that there are finitely many future choices in $\DChSet$. 
%
\begin{ass}\label{ass: jump_size}
The term $\vert f(x,\sigma^{\dVecVar}(x)) - f(x,\sigma^{\sVecVar}(x))\vert$ is bounded below by a constant $D$ for all $\dVecVar,\sVecVar\in \DChSet$, $\dVecVar\not = \sVecVar$ and $x\in S$. 
\end{ass}
%
Second, we assume that the rate of change of the next period states mapped by the transition and policy functions is uniformly bounded above by a common constant. Assumption \ref{ass: jump_size} is satisfied when the agent's policy function contains finitely many jumps. Assumption \ref{ass: lipshitz} is straightforward to satisfy in typical economic applications, where the policy function, conditional on a future stochastic sequence of discrete choices ($\sigma^{d}$), has a Lipshitz bound. This bound is commonly met since $\sigma^{d}$ is the solution to a concave optimization problem, in which the optimal policy function exhibits well-defined maximal and minimal marginal propensities to save \citep{Carroll2023}. 
%
% Note to Akshay: in the case with infinite dimensional state-space, this assumption should be weakened to that different policy functions that are present in the endogenous grid have a jump size bounded below. because in the above, if there is smoothing, the above will not hold. 

\begin{ass}\label{ass: lipshitz}
The family of functions $x\mapsto f(x,\sigma^{\dVecVar}(x))$ for $\dVecVar\in \DChSet$ have a common Lipshitz constant $M$ and an inverse Lipshitz constant $M_{\text{INV}}$.
\end{ass}

Finally, we assume each pair $Q^{\dVecVar}$ and $Q^{\sVecVar}$ either never intersect or crosses exactly once in $S$. A sufficient condition for this assumption to be satisfied is prudence, that is, $V^{\prime\prime\prime}(x)>0$. While more general conditions are also possible, standard applications in heterogenous agent macroeconomics and structural microeconometrics meet this requirement (\textcolor{red}{ref}). To state the assumption formally, for any distinct pair $(\dVecVar,\sVecVar)\in \mathbb{D}\times\mathbb{D}$ with $\dVecVar\neq \sVecVar$, define $\Delta_{\dVecVar,\sVecVar}(x) \colon = 
   Q^{\dVecVar}(x) - Q^{\sVecVar}(x)$. 

\begin{ass}\label{ass:single_crossing}
For every distinct pair $(\dVecVar,\sVecVar)$ in $\mathbb{D}\times\mathbb{D}$,
%
\begin{equation*}
   \Delta_{\dVecVar,\sVecVar}(x)
   \;>\; 0
   \quad\text{for all } x\in S,
   \qquad
   \text{or}
   \qquad
   \Delta_{\dVecVar,\sVecVar}(x)
   \;<\; 0
   \quad\text{for all } x\in S. 
\end{equation*}

\end{ass}

%%%%%%%%%%%%%%%%%%%

\begin{claim}\label{claim:grid_big}
 There exists $N$ such that for all $x_{j+1}^{\star}, x_{j}^{\star} \in \mathbb{X}^{\star}$ we have $\vert x_{j+1}^{\star}- x_{j}^{\star}\vert \leq \delta$ and $\frac{D}{\delta}> 2M$.
\end{claim}

\begin{claim}\label{claim:connected_optimal_regions}
Let $\DChSet$ be a finite set of discrete choices, and let $S$ be a compact state-space.
Suppose that for each $d\in \DChSet$, the region
\[
   S^{\dVecVar}
   \;\;=\;\;
   \bigl\{\,x \in S : \dVecVar = \arg\max_{\,\sVecVar \in \DChSet}\,G^{\sVecVar}(x)\bigr\}
\]
is a nonempty and connected subset of $S$.  Then there exists an integer $N^{\ast}>0$ such that for all exogenous grids of size $N \ge N^{\ast}$, each $S^{d}$ contains \emph{at least two distinct} grid points. 
\end{claim}


For each $d\in\DChSet$, define $\mathbb{X}^{d}\subset\{\left(\hat{x}_{i},\hat{x}_{i}^{\prime}\right)\}$ to be the points such that $Q^{d}(\hat{x}_{i}) = \hat{x}_{i}^{\prime}$. Let $\widetilde{Q}^{d}$ be the unique piecewise-linear function on the convex hull of $\mathbb{X}^{d}\cup \left\{0\right\}$ that interpolates $Q^{d}$. 
Set
\[
    \widetilde{U}(x) 
    \;=\;
    \max_{\,d\in\DChSet}\,\widetilde{Q}^{d}(x).
\]
Denote by $(\mathbb{X},\,\mathbb{X}^{\prime},\,\mathbb{Z},\,\mathbb{V})$ the output of the fast upper envelope scan (FUES) when run on $\{\hat{x}_{i}\}$ without forward or backward scans.

\begin{corollary}\label{corr:JDT}
There exists a jump detection threshold that satisfies $\frac{D}{\delta}-M \geq \bar{M} > M$.
\end{corollary}

For the following result, let $(\mathbb{X}, \mathbb{X}^{\prime}, \mathbb{Z}, \mathbb{V})$ be the tuple of outputs of Algorithm \ref{alg: NCE_scan} without the forward and backward scans.

\begin{proposition}\label{prop: main}
Let Assumptions \ref{ass: jump_size}-\ref{ass:single_crossing} hold. There exists grid size $N$ large enough such that there exists a jump detection threshold $\bar{M}$ with  $\frac{D}{\delta}-M \geq \bar{M} > M$ and
$\hat{v}_{i} = \widetilde{U}(\hat{x}_{i})$ for each $i$. 
\end{proposition}

\subsection{Proof of main result}

%% Todo AAS: Define $\tilde{S}^{d}$

We prove the Proposition \ref{prop: main} by showing

\begin{equation}\label{eq: claim_induct}
 \hat{x}_{i} \in \mathbb{X} \Longleftrightarrow  \hat{x}_{i} \in \widetilde{\mathbb{X}}^{\star}, \qquad \qquad \hat{x}_{i}\in \hat{\mathbb{X}},
\end{equation}

for all  $i\leq \vert \hat{\mathbb{X}} \vert$. In particular, we will show that if \eqref{eq: claim_induct} holds for all $i$ with $i\leq l$ for some $l\leq \vert \hat{\mathbb{X}} \vert$, then it will hold for all $i$ with $i\leq l+1$. Thus, by the principle of induction, the claim will hold for all $i$ with $i\leq l$ for any $l$, including $l = \vert \hat{\mathbb{X}} \vert$.

As such, fix $l$ with $l \leq \vert \hat{\mathbb{X}}\vert$ and make the inductive assumption that \eqref{eq: claim_induct} holds for all $i$ with $i\leq l$. Note that $\mathbb{X}$ is a subsequence of $\hat{\mathbb{X}}$ which we can denote as $\mathbb{X} = \{\hat{x}_{i_j}\}_{j=0}^{N_{F}}$, where $N_{F}$ is the number of points recovered by the FUES algorithm. As such, for any $k$ such that $i_{k}\leq l$, let $\{\hat{x}_{i_0}, \dots, \hat{x}_{i_j},\dots, \hat{x}_{i_{k}}\}$ denote the first $k$ points of $\mathbb{X}$.

\subsubsection{\textbf{Proof that $\hat{x}_{i} \in \mathbb{X}\Longrightarrow  \hat{x}_{i} \in \widetilde{\mathbb{X}}^{\star} $ holds for all $i$ with $i\leq l+1$.}}

For this part of the proof, fix $k$ such that $k = \max\left\{j:i_{j}\leq l\right\}$. If $i_{k+1}> l+1$, then the proof of this part is complete since $\hat{x}_{i} \in \mathbb{X} \Longrightarrow  \hat{x}_{i} \in \widetilde{\mathbb{X}}^{\star}$ holds for all $i< l+1$ by the inductive assumption. On the other hand, suppose  $i_{k+1}= l+1$. We wish to now show that $x_{i_{k+1}}\in \widetilde{\mathbb{X}}^{\star}$ (where the point $\hat{x}_{i_{k+1}}$ satisfies  $\hat{x}_{i_{k+1}}\in \mathbb{X}$). There are two cases. The first case arises if a right turn is made on the value correspondence, $\frac{\hat{v}_{i_{k+1}} - \hat{v}_{i_k}}{\hat{x}_{i_{k+1}}-\hat{x}_{i_k}} \leq \frac{\hat{v}_{i_{k}} - \hat{v}_{i_{k-1}}}{x_{i_k}-\hat{x}_{i_{k-1}}}$. The second case arises if a left turn is made on the value correspondence, $\frac{\hat{v}_{i_{k+1}} - \hat{v}_{i_k}}{\hat{x}_{i_{k+1}}-\hat{x}_{i_k}} > \frac{\hat{v}_{i_k} - \hat{v}_{i_{k-1}}}{\hat{x}_{i_k}-\hat{x}_{i_{k-1}}}$. 

%% AAS TO AAS: Why is the claim true for If $i_{k+1}> l+1$? Should we provide further details?
%% AAS to AAS: The progression of logic above is correct, however clarity could be improved. 

\textbf{[Part 1, Case 1: Right turn]}

Suppose $\hat{v}_{i_{k}} = Q^{\mVecVar}\bigl(\hat{x}_{i_{k}}\bigr)$ for some future sequence of discrete choices $\mVecVar$, and let $\hat{v}_{k+1} = Q^{\hat{\mVecVar}}\bigl(\hat{x}_{l+1}\bigr)$, where $\hat{\mVecVar}$ may or may not be the same as $\mVecVar$.  
We divide the analysis into two sub-cases.

\textit{Sub-case (i): The point $\hat{x}_{k+1}$ leads to the same future choice sequence as $\hat{x}_{k}$ and $\hat{\mVecVar} = \mVecVar$}

%% AAS TO ASS: WHAT IS \delta? make the link relating to the number of points in the exog.

By \textcolor{red}{Claim}~\ref{claim:grid_big}, the distance $\vert \hat{x}_{k+1} - \hat{x}_{i_{k}}\vert$ cannot exceed $\delta$.
If $\hat{x}_{k+1}\in \mathbb{X}^{\star}$, $\hat{x}_{k+1}$ lies before the next turning point, thus $\hat{x}_{k+1}$ must also belong to $\widetilde{\mathbb{X}}^{\star}$\textcolor{blue}{(detail in appx?)}.

%%AAS TO AAS: The why above has to do with whether we need further details on interpolation. 

On the other hand, if $\hat{x}_{k+1}\not \in \mathbb{X}^{\star}$, it must be that $\hat{x}_{k+1}\not \in S^{\mVecVar}$.  Thus, $\hat{x}_{k+1}$ is strictly greater than the smallest turning point greater than $\hat{x}_{k}$, which we denote by $x^{\text{TP}}$. Now, let $\hat{x}^{\dagger}> x^{\text{TP}}$ be the first optimal point in $\mathbb{X}^{\star}$ strictly greater than $x^{\text{TP}}$. The point $\hat{x}^{\dagger}$ leads the  optimal future choice sequence $\qVecVar \neq \mVecVar$.  
Recall that $\hat{x}_{k+1}' = f\!\bigl(\hat{x}_{l+1},\,\sigma_{\mVecVar}(\hat{x}_{l+1})\bigr)$ and $\hat{x}^{\prime,\dagger} =  f\!\bigl(\hat{x}^{\dagger},\,\sigma_{\qVecVar}(\hat{x}^{\dagger})\bigr)$.
%
% AAS to AAS: check the weak inequality above in terms of the main propostions statement. 
%
If $\hat{x}^{\dagger}$ is the minimal element in $\widetilde{S}^{\qVecVar}\cap \mathbb{X}$, let $\hat{x}^{\dagger +1}\widetilde{S}^{\qVecVar}\cap \mathbb{X}$ be the next grid point on the same future discrete choice sequence $\qVecVar$. Now, extrapolate the line segment from $\bigl(\hat{x}^{\dagger},\,Q^{\qVecVar}(\hat{x}^{\dagger})\bigr)$ to $\bigl(\hat{x}^{\dagger+1},\,Q^{\qVecVar}(\hat{x}^{\dagger+1})\bigr)$ 
to the point $\hat{x}_{i_{k+1}}< \hat{x}^{\dagger}$ to form the evaluation of $\widetilde{Q}^{\qVecVar}(\hat{x}_{i_{k+1}})$. However, the backward scan procedure eliminates $\hat{x}_{i_{k+1}}$ from the subsequence $\mathbb{X}$ if $\widetilde{Q}^{\qVecVar}(\hat{x}_{i_{k+1}})> \widetilde{Q}^{\mVecVar}(\hat{x}_{i_{k+1}}) = \hat{v}_{i_{k+1}}$. This implies that $\widetilde{Q}^{\qVecVar}(\hat{x}_{i_{k+1}})\leq  \widetilde{Q}^{\mVecVar}(\hat{x}_{i_{k+1}}) = \hat{v}_{i_{k+1}}$ and thus $\hat{x}_{i_{k+1}}\in \widetilde{\mathbb{X}}^{\star}$. 

%%  AAS TO AAS, WHY Will the backward scan do this? The backward scan will detect jumps in all the points before x_{ik} and the point x^{\dagger} and hence extrapolate. '
%% We have to clarify above that x^{|dagger} is indeed the first optimal point that is picked up. 

If $\hat{x}^{\dagger}$ is not the minimal element in $\widetilde{S}^{\qVecVar}\cap \mathbb{X}$, then let $\hat{x}^{\dagger -1}\widetilde{S}^{\qVecVar}\cap \mathbb{X}$ be the grid point on the same future discrete choice sequence $\qVecVar$ to the left of $\hat{x}^{\dagger}$. Note that this point must be to the left of $x^{\text{TP}}$. Thus, the line segment from $\bigl(\hat{x}^{\dagger},\,Q^{\qVecVar}(\hat{x}^{\dagger})\bigr)$ to $\bigl(\hat{x}^{\dagger-1},\,Q^{\qVecVar}(\hat{x}^{\dagger-1})\bigr)$ can be evaluated at the point $\hat{x}_{i_{k+1}}< \hat{x}^{\dagger}$ to  evaluate $\widetilde{Q}^{\qVecVar}(\hat{x}_{i_{k+1}})$. Hence, by the backward scan, \textcolor{blue}{(ddetail in appx?)} if  $\widetilde{Q}^{\qVecVar}(\hat{x}_{i_{k+1}})$ evaluates strictly above $\hat{v}_{i_{k+1}}$, then $\hat{x}_{i_{k+1}}$ is removed as suboptimal.  Otherwise, if $\widetilde{Q}^{\qVecVar}(\hat{x}_{i_{k+1}})$ is below $\hat{v}_{i_{k+1}}$, the point remains in $\hat{x}_{i_{k+1}}\in \mathbb{X}$ and must lie on the envelope, so $\hat{x}_{i_{k+1}} \in \widetilde{\mathbb{X}}^{\star}$.

\textit{Sub-case (ii): The point $\hat{x}_{i_{k+1}}$ switches to a distinct future sequence $\hat{\mVecVar}\neq \mVecVar$.}

If $\qVecVar\neq \mVecVar$ and $\hat{x}_{i_{k+1}}\leq x^{\text{TP}}$, then by Claim~\ref{clm: jump_detect}, the FUES algorithm detects a discontinuous jump and removes $\hat{x}_{i_{k+1}}$.  Hence we must have $\hat{x}_{i_{k+1}} > x^{\text{TP}}$.  In this case the FUES algorithm also detects a discontinuous jump, and if $\hat{x}_{i_{k+1}}$ remains in $\mathbb{X}$, the forward scan must not have eliminated it.  Specifically, the forward scan looks for a point $\hat{x}^{\dagger} > x^{\text{TP}}$ on the same future discrete choice sequence path $\mVecVar$.  
If such a $\hat{x}^{\dagger}$ exists, the algorithm compares the slope from $(\hat{x}_{i_{k}},\hat{v}_{i_{k}})$ to $(\hat{x}^{\dagger},Q^{\mVecVar}(\hat{x}^{\dagger}))$ with the slope from $(\hat{x}_{i_{k}},\hat{v}_{i_{k}})$ to $(\hat{x}_{i_{k+1}},\hat{v}_{i_{k+1}})$. If the former slope is strictly greater, then $\hat{x}_{i_{k+1}}$ is deemed suboptimal and removed. If the slope from $(\hat{x}_{i_{k}},\hat{v}_{i_{k}})$ to $(\hat{x}^{\dagger},Q^{\mVecVar}(\hat{x}^{\dagger}))$ is less than the slope from $(\hat{x}_{i_{k}},\hat{v}_{i_{k}})$ to $(\hat{x}_{i_{k+1}},\hat{v}_{i_{k+1}})$, then $\hat{v}_{i_{k+1}})= Q^{\qVecVar}(\hat{x}^{\dagger})$. If no such $\hat{x}^{\dagger}$ is found, $\hat{x}_{i_{k+1}}$ is not an element of $\mathbb{X}$

%% AS TO AS: think again about what happens when If no such $\hat{x}^{\dagger}$ is found ...does this still generate the upper envelope of the interpoalnts?
% AAS TO AAS: The above seems alittle loose. 

\textbf{[Part 1, Case 2: Left turn]}

Let $\frac{\hat{v}_{i_{k+1}} - \hat{v}_{i_{k}}}{\hat{x}_{i_{k+1}}-\hat{x}_{i_k}} > \frac{\hat{v}_{i_k} - \hat{v}_{i_{k-1}}}{\hat{x}_{i_k}-\hat{x}_{i_{k-1}}}$. Then by Claim \ref{clm: increasing_g} in the Appendix, we must have that $U(\hat{x}_{i_{k+1}},\hat{z}_{i_{k+1}}) = V(\hat{x}_{i_{k+1}})$, and so $\hat{x}_{i_{k+1}}\in \mathbb{X}^{\star}$. We will now show that this implies $\hat{x}_{i_{k+1}}\in \widetilde{\mathbb{X}}^{\star}$. Suppose for contradiction that $\hat{x}_{i_{k+1}}\not \in \widetilde{\mathbb{X}}^{\star}$. Then by definition of $\widetilde{\mathbb{X}}^{\star}$, there must exist a point $\hat{x}^{\dagger} \in \widetilde{\mathbb{X}}^{\star}$ such that $\hat{x}^{\dagger} < \hat{x}_{i_{k+1}}$ and a point $\hat{x}^{\dagger +1} \in \widetilde{\mathbb{X}}^{\star}$ such that $\hat{x}^{\dagger +1} > \hat{x}_{i_{k+1}}$. Moreover, we must have that $\hat{x}^{\dagger}$ and $\hat{x}^{\dagger +1}$ lie on the same future choice sequence $\rVecVar$. But since each $Q^{\rVecVar}$ is non-decreasing in $\rVecVar$, we must have that $Q^{\hat{\mVecVar}}(\hat{x}_{i_{k+1}}) < Q^{\rVecVar}(\hat{x}_{i_{k+1}})$. This is a contradiction to $\hat{x}_{i_{k+1}}\in \mathbb{X}^{\star}$, and hence $\hat{x}_{i_{k+1}}\in \widetilde{\mathbb{X}}^{\star}$.

%% TODO AAS: Check the above carefully. 

\subsubsection{Proof that $\hat{x}_{i} \in \mathbb{X} \Longleftarrow  \hat{x}_{i} \in \widetilde{\mathbb{X}}^{\star} $ holds for $i$ with $i\leq l+1$.}

Now we show that if a candidate point $\hat{x}_{l+1}$, with $\hat{x}_{l+1}\in \widetilde{\mathbb{X}}^{\star}$ and $ V(\hat{x}_{l+1}) = Q(\hat{x}_{l+1},\hat{z}_{l+1})$, then $\hat{x}_{l+1}\in \mathbb{X}$. Consider the first case of a right turn. 

\textbf{[Part 2, Case I: Right turn]}

Suppose $\hat{x}_{l+1}$ generates a right turn with respect to the line segment joining $\hat{x}_{i_{k-1}}$ and $\hat{x}_{i_{k}}$. Formally, let
%
\begin{equation}\label{eq:rt_slope}
    \frac{\hat{v}_{l+1} - \hat{v}_{i_{k}}}{\hat{x}_{l+1} - \hat{x}_{i_{k}}}
    \;\;\le\;\;
    \frac{\hat{v}_{i_{k}} - \hat{v}_{i_{k-1}}}{\hat{x}_{i_{k}}-\hat{x}_{i_{k-1}}},
\end{equation}
%
where $\hat{v}_{i_{k}} = Q^{\mVecVar}\bigl(\hat{x}_{i_{k}}\bigr)$ for some future choice sequence $\mVecVar\in\mathbb{D}$.  In addition, let $\hat{v}_{l+1} = Q^{\qVecVar}\bigl(\hat{x}_{l+1}\bigr)$ for $\qVecVar\in \mathbb{D}$.  The aim is to demonstrate that $\hat{x}_{l+1} \in \mathbb{X}$. There are two sub-cases to examine, $\qVecVar = \mVecVar$ and $\qVecVar \neq \mVecVar$.

\paragraph{Sub-case (i): $\qVecVar = \mVecVar$} Since $\hat{x}_{l+1}$ and $\hat{x}_{i_{k}}$ lie at most $\delta$ apart by Claim~\ref{claim:grid_big}, we now show that no policy jump can occur if both points pertain to the same discrete choice $m$.  To see why, denote
%W
\begin{equation*}
   \hat{x}_{l+1}' 
   \;=\; 
   f\Bigl(\hat{x}_{l+1},\,\sigma_{\mVecVar}\bigl(\hat{x}_{l+1}\bigr)\Bigr),
   \quad
   \hat{x}_{i_{k}}' 
   \;=\; 
   f\Bigl(\hat{x}_{i_{k}},\,\sigma_{\mVecVar}\bigl(\hat{x}_{i_{k}}\bigr)\Bigr).
\end{equation*}
%
Under Assumption~\ref{ass: jump_size}, it must hold that 
%
\begin{equation}
   \Bigl|\hat{x}_{l+1}' - \hat{x}_{i_{k}}'\Bigr|
   \;\le\; 
   M\,\Bigl|\hat{x}_{l+1} - \hat{x}_{i_{k}}\Bigr|,
\end{equation}
%
which then implies
%
\begin{equation}
   \frac{\Bigl|\hat{x}_{l+1}' - \hat{x}_{i_{k}}'\Bigr|}
        {\Bigl|\hat{x}_{l+1} - \hat{x}_{i_{k}}\Bigr|}
   \;\le\; 
   M< \frac{D}{\delta}-M, 
\end{equation}
%
where the second inequality follows from Corollary~\ref{corr:JDT}.

On the other hand, the jump-detection threshold in line~\ref{lst:line: right_no_jump} of Algorithm~\ref{alg: NCE_scan} requires that if 
%W
\begin{equation}
   \frac{\Bigl|\hat{x}_{l+1}' - \hat{x}_{i_{k}}'\Bigr|}
        {\Bigl|\hat{x}_{l+1} - \hat{x}_{i_{k}}\Bigr|}
   \;<\;
   \frac{D}{\delta} - M,
\end{equation}
%W
then FUES does remove the point $\hat{x}_{l+1}$ from $\mathbb{X}$. Thus, $\hat{x}_{l+1}\in \mathbb{X}$ as was to be shown. 

\paragraph{Sub-case (ii):$\qVecVar \neq \mVecVar$}

In this sub-case, $\hat{x}_{i_{k-1}}$ and $\hat{x}_{i_{k}}$ both lie in the region of the state-space that yields the optimal future choice sequence $\mVecVar$, but $\hat{x}_{l+1}$ satisfies $V(\hat{x}_{l+1}) = Q^{\qVecVar}\bigl(\hat{x}_{l+1}\bigr)$ for some $\qVecVar\neq \mVecVar$.  
By the forward scan procedure, one observes that $x_{l+1}$ must be strictly below $\widetilde{Q}^{\mVecVar}$ from $\hat{x}_{i_{k-1}}$ to $\hat{x}_{i_{k}}$ (\textcolor{red}{detail in appx?}). 

%% Below is a poor explanation of the why above, needs refinement 
%To see this, note that if both $\hat{x}_{i_{k-1}}$ and $\hat{x}_{i_{k}}$ lie in $m$'s region, then
%\[
%   \widetilde{Q}^{m}(x)
%   \;=\;
%   \widetilde{Q}^{m}\Bigl(\hat{x}_{i_{k-1}}
%          + \theta\bigl(x - \hat{x}_{i_{k-1}}\bigr)\Bigr)
%\]
%applies for all $x\in[\hat{x}_{i_{k-1}},\hat{x}_{i_{k}}]$ and by extrapolation outside $[\hat{x}_{i_{k-1}},\hat{x}_{i_{k}}]$ as well, given the linear extension in $\widetilde{Q}^{m}$.  
%Hence the slope test in \eqref{eq:rt_slope} ensures $\hat{x}_{l+1}$ falls below that linear extension of $\widetilde{Q}^{m}$, implying $Q^{m}(\hat{x}_{l+1}) < \hat{v}_{l+1} = Q^{\qVecVar}(\hat{x}_{l+1})$.  
%Thus $\hat{x}_{l+1}$ cannot reside on $\widetilde{Q}^{m}$, leaving $\widetilde{Q}^{\qVecVar}$ as the unique interpolant containing $\hat{x}_{l+1}$; consequently $\hat{x}_{l+1}\in \widetilde{\mathbb{X}}^{\star}$.

%WA final possibility arises if there exists a point $\hat{x}^{\prime,\dagger} \in Z^{\mVecVar}$ with $\hat{x}^{\prime,\dagger} > \hat{x}_{l+1}'$, where $\hat{x}_{i_{k}}'=\max Z^{\mVecVar}$ no longer holds.  
%In that event, the line from $\bigl(\hat{x}_{i_{k}},\hat{v}_{i_{k}}\bigr)$ to $\bigl(\hat{x}^{\dagger},Q^{m}(\hat{x}^{\dagger})\bigr)$ may dominate or fall below $\bigl(\hat{x}_{l+1},\hat{v}_{l+1}\bigr)$.  
%By Assumption~\ref{ass: grid_big}, the distance between $\hat{x}^{\dagger}$ and $\hat{x}_{l+1}$ can be bounded by $\delta$, so $\bigl|\hat{x}^{\dagger}-\hat{x}_{l+1}\bigr|
%   \;\leq\;
%   \delta$. 

%Algorithm~\ref{alg: NCE_scan} (the forward-scan step) compares the slope from $\hat{x}_{i_{k}}$ to $\hat{x}^{\dagger}$ with that from $\hat{x}_{i_{k}}$ to $\hat{x}_{l+1}$.  
%If $\hat{x}_{l+1}$ exceeds the secant between $(\hat{x}_{i_{k}},\hat{v}_{i_{k}})$ and $(\hat{x}^{\dagger},\,Q^{m}(\hat{x}^{\dagger}))$, then line~\ref{lst:line: left_turn} of the algorithm confirms an upward turn, so the point $\hat{x}_{l+1}$ is recognized as optimal on $\widetilde{Q}^{\qVecVar}$.  
%If $\hat{x}_{l+1}$ remains below that secant, the slope test excludes it.  
%Either way, each truly optimal point (in particular, any point with $V(\hat{x}_{l+1})=Q^{\qVecVar}$) is \emph{retained} if and only if it lies on the envelope.  
%Hence $\hat{x}_{l+1}\in \widetilde{\mathbb{X}}^{\star}$, completing the argument under $\qVecVar\neq m$.

%In summary, whether $\qVecVar=m$ or $\qVecVar\neq m$, if $\hat{x}_{l+1}\in \mathbb{X}$ satisfies the right-turn slope \eqref{eq:rt_slope}, then $\hat{x}_{l+1}\in \widetilde{\mathbb{X}}^{\star}$.  
%This concludes Part~2, Case~1.

\textbf{[Part 2: Left turn]}

Let us now assume that
%
\begin{equation}\label{eq:lt_slope}
    \frac{\hat{v}_{l+1} - \hat{v}_{i_{k}}}{\hat{x}_{l+1} - \hat{x}_{i_{k}}}
    \;>\;
    \frac{\hat{v}_{i_{k}} - \hat{v}_{i_{k-1}}}{\hat{x}_{i_{k}}-\hat{x}_{i_{k-1}}}.
\end{equation}
%
Under this condition, a standard slope comparison indicates that the point $\hat{x}_{l+1}$ produces a strictly \emph{higher} secant slope than the line segment from $(\hat{x}_{i_{k-1}}, \hat{v}_{i_{k-1}})$ to $(\hat{x}_{i_{k}}, \hat{v}_{i_{k}})$.  
By line~16 of Algorithm~\ref{alg: NCE_scan}, the procedure interprets this higher slope as an upward or ``left-turn'' increase in the value function.  
Consequently, $\hat{x}_{l+1}$ is retained and $\hat{x}_{l+1}\in \mathbb{X}$.


%%%%%%%%%%%%%%%%%%%%%%%%%%%%%%%%%%%%%%%%%%%%%

\section{Conclusions}

This paper develops a fast upper envelope scan (FUES) method to solve any discrete-continuous stochastic dynamic optimization models where the optimization problem has a one-dimensional representation. FUES works by efficiently computing the upper envelope of the value correspondence generated by the EGM, overcoming the limitations of current methods that fail in applications involving non-monotonic policy functions. We demonstrate the wide applicability and massive computational advantages of FUES over competing upper envelope methods like DC-EGM \citep{Iskhakov2017} and RFC \citep{DS2024RFC}, as well as over hybrid methods that use numerical optimization like NEGM \citep{Druedahl2021} via three workhorse applications in the literature. Specifically, we find FUES to achieve computational speed gains of generally an order of magnitude above current methods, while maintaining solution accuracy. Moreover, in applications where DC-EGM cannot be applied due to non-monotonicity of the policy function, FUES also leads to up to an order of magnitude better Euler error over NEGM. Crucially, FUES scales well with model complexity, providing consistently high performance even when the endogenous grids become more irregular due to uncertainty or non-monotonicity.

One final remark relates to FUES being ultimately designed for models where the EGM grids of the potential solution points can be collapsed to one-dimensional ones, as it relies on the sequential ordering of points in the endogenous grid to perform the scan. This approach is not yet feasible in higher dimensions, where endogenous grids are irregular and lack a natural order for scanning. A promising area for future research is thus extending FUES to multiple dimensions by incorporating  techniques such as Hilbert curves or other space-filling methods to produce an ordering of irregular multidimensional points. 

%%%%%%%%%%%%%%%%%%%%%%%%%%%%%%%%%%%%%%%%%%%%%

\bibliographystyle{apalike}
\bibliography{DCegm} 

%%%%%%%%%%%%%%%%%%%%%%%%%%%%%%%%%%%%%%%%%%%%%

\newpage 

\appendix

{\huge Online Appendix}

%%%%%%%%%%%%%%%%%%%

\section{Existing methods}\label{appx:existing}

\paragraph{Monotone Segment Selection \citep{Iskhakov2017}}

To illustrate the approach by \cite{Iskhakov2017}, refer again to the right panel of Figure~\ref{fig: FUES_1}. In that panel, note that the set of points
%
\[
\bigl\{(\hat{x}_{1}, \hat{x}_{1}^{\prime}),\, (\hat{x}_{2}, \hat{x}_{2}^{\prime}),\, (\hat{x}_{3}, \hat{x}_{3}^{\prime}),\, (\hat{x}_{4}, \hat{x}_{4}^{\prime}),\, (\hat{x}_{7}, \hat{x}_{7}^{\prime}),\, (\hat{x}_{9}, \hat{x}_{9}^{\prime})\bigr\}
\]
forms a monotone segment when plotted along the vertical axis. When these points are ordered by the exogenous variable $\hat{x}^{\prime}$ on the exogenous grid, they appear consecutively and are associated with a single future choice-specific policy function. In contrast, the remaining points 
\[
\bigl\{(\hat{x}_{5}, \hat{x}_{5}^{\prime}),\, (\hat{x}_{6}, \hat{x}_{6}^{\prime}),\, (\hat{x}_{8}, \hat{x}_{8}^{\prime})\bigr\}
\]
belong to a different future choice-specific policy function, which is also monotone over its own domain. DC-EGM exploits the monotonicity of the policy function by performing the following steps:
%
\begin{enumerate}[label=(\roman*)]
    \item Identify distinct monotone segments within the exogenous grid.
    \item Create an interpolant using the current period values over these segments separately.
    \item Locate the crossing points between these interpolated segments (as illustrated in the left panel of Figure~\ref{fig: FUES_1}).
    \item Remove candidate points that lie below the dominating upper envelope on either side of each crossing point.
\end{enumerate}

%Moreover, because the segmentation is performed on the exogenous grid, the number of segments cannot be fixed a priori, which prevents efficient vectorized computation of the upper envelope algorithm. 

\paragraph{Local Triangulation Method \citep{Druedahl2017}}

\cite{Druedahl2017} propose a method that, in multiple dimensions, constructs a local triangulation over the exogenous grid and then compares the resulting piecewise interpolated upper envelope on the endogenous grid. In a single-dimensional setting, their procedure reduces to a \emph{pairwise segment comparison} on an interval of the endogenous wealth space. Concretely, let 
\(\{\hat{x}_{i}, \hat{x}_{i}^{\prime}, \hat{v}_{i}\}_{i=1,\dots,N}\)
be exogenous- and endogenous-grid pairs as in Figure~\ref{fig: FUES_1}, ordered by the $\hat{x}_{i}^{\prime}$ in the exogenous grid. Also fix a uniform grid of current assets $\mathbb{X}_{i}^{U}$. The algorithm proceeds as follows:

\begin{enumerate}[label=(\roman*)]
    \item Identify all consecutive pairs of exogenous points 
    \(\bigl(\hat{x}_{i}^{\prime}, \hat{x}_{i+1}^{\prime}\bigr)\) in the exogenous grid.
    
    \item For each pair \(\bigl(\hat{x}_{i}^{\prime}, \hat{x}_{i+1}^{\prime}\bigr)\), define a bounding interval 
    \(\bigl[\underline{x}_{i},\overline{x}_{i}\bigr]\)
    for current assets \(x\) in the endogenous space. %% AS TO AS do we need to say EGM used?
    
    \item For each bounding interval \(\bigl[\underline{x}_{i},\overline{x}_{i}\bigr]\), select those points in \(\mathbb{X}^{U}\) (a uniform grid of current assets) that lie within the interval. Denote this subset by \(\mathbb{X}_{i}^{U}\).

    \item For each point \(\hat{x}^{U}_{j}\in \mathbb{X}_{i}^{U}\), interpolate the value 
    \(v_{i,j}^{U}\)
    using data from 
    \(\bigl(\hat{x}_{i}^{\prime}, \hat{v}_{i}\bigr)\)
    and
    \(\bigl(\hat{x}_{i+1}^{\prime}, \hat{v}_{i+1}\bigr)\).
    If multiple bounding intervals overlap at \(\hat{x}^{U}_{j}\), compare the interpolated values \(v_{i,j}^{U}\) across all such intervals \(i\), and select the interval \(i^{\star}\) yielding the highest value.

    \item The pair 
    \(\bigl(\mathbb{X}_{i}^{U}, \mathbb{V}_{i}^{U}\bigr)\)
    serves as an interpolant for approximating the value function. Similarly, an approximate policy function can be constructed by interpolating the optimal policy  associated with each $v_{i,j}^{U}$. %% CAN this be clearer?
\end{enumerate}

\section{Applications}

\subsection{Retirement coice -- shocks}\label{appx:retShocks}
To further illustrate the ability of FUES to robustly handle more complex endogenous grids with non-monotone segments and isolated points, we also compare the performance of FUES and DC-EGM in a model where an extreme value taste shock with scale parameter $\bar{s}$ is added to the utility of workers (see \cite{Iskhakov2017} for details of the model with taste shocks). The resulting endogenous grid is illustrated in Figure \ref{fig:dcegm_sigma1}, where decreasing non-monotone parts of the endogenous grid lead to a significant number of additional segments the DC-EGM algorithm has to process. The presence of these smaller segments leads to slower performance for DC-EGM (see Table \ref{table:RETREST2}), with FUES now close to an order of magnitude faster than DC-EGM, while still maintaining its higher accuracy. 

%\FloatBarrier
\begin{figure}[H]
\centering
\includegraphics[scale=0.65]{plots/ret_vf_aprime_all_3000_cf_17_sigma_035.png}
\caption{Non-monotone endogenous grid points in presence of taste shocks.}\label{fig:dcegm_sigma1}
\end{figure}

% Are we using subscripts for endogenous grid points and time?

\section{Endogenous grid-specific jump detection}\label{sec:ENDJTAppx}

In practice, we find that selecting a reasonable constant value for $\bar{M}$ through experimentation (and subsequently verifying the upper envelope is recovered) works well. However, an endogenous jump detection threshold $\bar{M}^{\star}_{i}$ can also be theoretically derived based on the curvature of the policy functions, and easily computed for each grid point. To this effect, note that the role of the jump detection threshold is to detect a point $\hat{x}_{i+1}$ that does not lie on the same future choice-specific policy function as $\hat{x}_{i}$. Consider now that for any real-valued continuously differentiable function $f$ on an interval $I$, by the Mean Value Theorem, we must have that $\vert f(y)- f(x)\vert \leq L\vert y-x\vert$, where $L = \sup\{|f'(z)| : z \in I\}$. Let $\sigma_{t}^{\textbf{d}}$ be the  asset policy function at time $t$, conditional on a future sequence of discrete choices $\textbf{d}$ (recall the discussion below \eqref{eq: Bell_Q}). The functions $a \mapsto \sigma^{\textbf{d}}_{t}(a)$ are continuously differentiable and convex \citep{Carroll2004} since all future discrete choices are held fixed. For any exogenous grid point $\hat{x}_{i+1}$ to imply the same future discrete choices as $\hat{x}_{i}$, we must then have that
%
\begin{align}\label{eq:endbound}
\vert \sigma^{\textbf{d}_{i}}_{t}(\hat{x}_{i}) - \sigma^{\textbf{d}_{i+1}}_{t}(\hat{x}_{i+1}) \vert 
&\leq \bar{M}^{\star}_{i}\vert \hat{x}_{i+1} - \hat{x}_i\vert ,
\end{align}
%
where $\bar{M}^{\star}_{i} = \max\left\{\left| \partial \sigma^{\textbf{d}_{i+1}}(\hat{x}_{i+1})\right|, \left| \partial \sigma^{\textbf{d}_{i}}(\hat{x}_i) \right|\right\}$ and we use $\partial$ to denote the derivative of the policy function with respect to financial assets. The term $\textbf{d}_{i}$ above is the future sequence of discrete choices implied by the exogenous grid point $i$. If two points do not satisfy the above inequality, then they cannot lie on the same future choice-specific policy function and we have `detected a jump` in the policy function. 

Starting from the exogenous grid points, we can use the Euler equation for the problem to derive the terms $\partial \sigma^{\textbf{d}_{i}}(\hat{x}_{i})$ analytically and recursively as follows:
%
\begin{equation}
\partial \sigma^{\textbf{d}_{i}}_{t}(\hat{x}_{i}) =
\frac{(1+r)u^{\prime\prime}(\hat{c}_{i})}{u^{\prime\prime}(\hat{c}_{i}) + \beta(1+r)u^{\prime\prime}(\hat{c}_{i}^{\prime})[(1+r) - \partial \sigma^{\textbf{d}_{i}}_{t+1}(\hat{x}_{i}^{\prime}) ]},
\end{equation}
%
where $\hat{c}_{i}^{\prime}$ is the next period consumption implied by the exogenous grid point $i$. It is important to note that the dependence of the policy function derivatives $\partial \sigma^{\textbf{d}_{i}}_{t}$ on the future sequence of discrete choices is \textit{given implicitly} by the exogenous grid point $i$. Thus, the above derivative can be calculated by evaluating the RHS of the equation above as follows: (i) evaluate $u^{\prime\prime}(\hat{c}_{i})$ using the $\hat{c}_{i}$ implied by the standard EGM step, and then (ii) evaluate the next period value $u^{\prime\prime}(\hat{c}_{i}^{\prime})\left[ (1+r) - \partial \sigma^{\textbf{d}_{i}}_{t+1}(\hat{x}_{i}^{\prime})\right]$ using an interpolant of the function $u^{\prime\prime}(c^{\prime})[ (1+r) - \partial \sigma^{\textbf{d}}_{t+1}(x^{\prime})]$. In what follows, at a time $t$ iteration, the function $u^{\prime\prime}(c^{\prime})\left[ (1+r) - \partial \sigma^{\textbf{d}}_{t+1}(x^{\prime})\right]$
is evaluated on the exogenous grid points and stored in a grid $\hat{\mathbb{O}}_{t}^{\prime}$. (For period $T$, we can use 
$\partial \sigma^{\textbf{d}}_{T} = 0$ since all wealth is consumed.) In particular, fixing a time $t$ and noting $\hat{o}_{i}^{\prime}\in \hat{\mathbb{O}}^{\prime}_{t}$, define:
%
\begin{equation}
\hat{\Delta}_{i} = \frac{(1+r)u^{\prime\prime}(\hat{c}_{i})}{u^{\prime\prime}(\hat{c}_{i}) + \beta(1+r)\hat{o}_{i}^{\prime}} \approx \partial \sigma_{t}^{\textbf{d}_{i}}(\hat{x}_i),
\end{equation}
%
\begin{equation}
\hat{o}_{i} = u^{\prime\prime}(\hat{c}_{i})\left[(1+r) - \hat{\Delta}_{i} \right].
\end{equation}

Using the above, $\hat{\mathbb{O}}_{t}$ can be constructed from the points $\hat{o}_{i}$. The modified pseudo-code for FUES with endogenous jump detection becomes:\footnote{The FUES implementation with endogenous jump detection is provided as an additional option \href{http://github.com/akshayshanker/FUES_EGM}{\textcolor{blue}{here}}.}

\begin{tcolorbox}[
colframe=black!100,
colback=white!10,
coltitle=white!100!black,  
fonttitle=\bfseries,
adjusted title= Box 2: FUES method with endogenous jump detection]
\begin{enumerate}
    \item Compute $\hat{\mathbb{X}}_{t}$, $\hat{\mathbb{C}}_{t}$, $\hat{\mathbb{V}}_{t}$, $\hat{\mathbb{O}}_{t}$, $\hat{\mathbb{X}}_{t}^{\prime}$, $\hat{\mathbb{C}}^{\prime}$ and $\hat{\mathbb{O}}_{t}^{\prime}$  using standard EGM.
    \item Sort all sequences in order of the \textit{endogenous} grid $\hat{\mathbb{X}}_{t}$.
    \item Start from point $i= 2$.
    \item Compute $g_{i} = \frac{\hat{v}_{i} - \hat{v}_{i-1}}{\hat{x}_{i}- \hat{x}_{i-1}}$, $g_{i+1} = \frac{\hat{v}_{i+1} - \hat{v}_{i}}{\hat{x}_{i+1}- \hat{x}_{i}}$, $\hat{\Delta}_{i}$, and $\hat{\Delta}_{i+1} $.
    \item If the endogenous bound in \eqref{eq:endbound} is violated and a right turn is made ($g_{i+1}< g_{i}$), then remove point $i+1$ from grids $\hat{\mathbb{X}}_{t}$, $\hat{\mathbb{V}}_{t}$, $\hat{\mathbb{O}}_{t}$, $\hat{\mathbb{X}}_{t}^{\prime}$ and $\hat{\mathbb{O}}_{t}^{\prime}$. Otherwise, set $i= i+1$. 
    \item If $i+1\leq \vert \hat{\mathbb{X}}_{t}\vert$, then repeat from step 5. 
\end{enumerate}
\end{tcolorbox}

Once the scan is complete for time $t$, the refined grid $\mathbb{O}_{t}$ can be interpolated on the refined endogenous grid along with the refined value and policy function. We also note that if the policy functions are Lipshitz and there are only finitely many jumps, then the grid-specific bounds will be bounded and we will have effectively chosen a bound at each grid point such that the jump detection threshold at Assumption \ref{ass: JDT} in Section \ref{sec: theory} is satisfied. Hence, for a grid size large enough, FUES will eliminate all sub-optimal points without error. 

%%%%%%%%%%%%%%%%%

\section{Intermediate proofs}

The first intermediate result says that if a point makes a left turn from an optimal point, then it is also optimal. 

\begin{claim}\label{clm: increasing_g}
Fix the triple $\hat{x}_{i}, \hat{x}_{i+1}, \hat{x}_{i+2}$ for some $i$ and assume $\hat{x}_{i}, \hat{x}_{i+1}\in \mathbb{X}^{\star}$. If we have that

\begin{equation}\label{eq: left_turp_tp1}
    \frac{\hat{v}_{i+1}-\hat{v}_{i}}{\hat{x}_{i+1}- \hat{x}_{i}}< \frac{\hat{v}_{i+2}- \hat{v}_{i+1}}{\hat{x}_{i+2}- \hat{x}_{i+1}},
\end{equation}
%
then $\hat{x}_{i+2}\in \mathbb{X}^{\star}$.
\end{claim}

\begin{proof}
Fix $l$, with $l\in \DChSet$ and such that $V(\hat{x}_{i+1}) = G^{l}(\hat{x}_{i+1}, \sigma^{l}(\hat{x}_{i+1}))$. Suppose by contradiction that $\hat{x}_{i+2}\not\in \mathbb{X}^{\star}$. Suppose first that $V(\hat{x}_{i}) =  G^{l}(\hat{x}_{i}, \sigma^{l}(\hat{x}_{i}))$. Since $\hat{x}_{i+2}\not\in \mathbb{X}^{\star}$, we must have $\hat{v}_{i+2} =  G^{m}(\hat{x}_{i+2}, \sigma^{m}(\hat{x}_{i+2}))$ and $m\not = l$. By concavity and \eqref{eq: left_turp_tp1}, we must have $G^{m}(\hat{x}_{i+2}, \sigma^{m}(\hat{x}_{i+2}))> G^{l}(\hat{x}_{i+2}, \sigma^{l}(\hat{x}_{i+2}))$. This implies that $V(\hat{x}_{i+2}) = G^{p}(\hat{x}_{i+2}, \sigma^{p}(\hat{x}_{i+2}))$ for $p\not = l$ and $p\not = m$. Moreover, the functions $G^{p}$ and $G^{l}$ must cross at some point $\tilde{x}_{i}\in T_{p}$, with $\tilde{x}_{i} \in [\hat{x}_{i+1},\hat{x}_{i+2})$. However, we have now violated Assumption \ref{ass: intersection_opt} - item 1., since $\hat{x}_{i+2}$ is the first point after a crossing point and must be optimal. 

Now suppose $V(\hat{x}_{i}) =  G^{b}(\hat{x}_{i}, \sigma^{b}(\hat{x}_{i}))$ for some $b\not = l$. This implies a crossing point $\tilde{x}$ satisfies $\tilde{x}\in (\hat{x}_{i},\hat{x}_{i+1})$. However, $\tilde{x}\in (\hat{x}_{i},\hat{x}_{i+1})$ also yields a contradiction since by  Assumption \ref{ass: turning},  we must have $\tilde{x}\in \mathbb{X}^{\star}$. 
%and another crossing point $\h%t{\tilde{x}}\in [\hat{x}_{i+1},\hat{x}_{i+2}]$. By Assumption \ref{ass: turning} - item 2., we again have a contradiction since the assumption requires $v^{\star}_{j+2} = G^{l}(x^{\star}_{j+2}, \sigma^{l}(x^{\star}_{j+2}))$ or $v^{\star}_{j+2} = G^{p}(x^{\star}_{j+2}, \sigma^{p}(x^{\star}_{j+2}))$, where $p$ is such that $V(\hat{x}_{i+1}) = G^{p}(\hat{x}_{i+1}, \sigma^{p}(\hat{x}_{i+1}))$. Thus, in the case where $V(\hat{x}_{i}) \not =  G^{l}(\hat{x}_{i}, \sigma^{l}(\hat{x}_{i}))$, we must have $\hat{x}_{i+2} \in \mathbb{X}^{\star}$, completing the proof. 
\end{proof}

% Proof above could be written more clearly with more detail. 


\begin{claim}\label{clm: x_k_not_TP_1}
Consider the setting of Part 1 of the proof of Proposition \ref{prop: main}. Let $\tilde{x}$ be the smallest value in $T_{p}$ such that $\tilde{x}> \hat{x}_{i_{k}}$. If $\frac{\hat{v}_{i_{k+1}} - \hat{v}_{i_k}}{\hat{x}_{i_{k+1}} -\hat{x}_{i_k}} \leq \frac{\hat{v}_{i_k} - \hat{v}_{i_{k-1}}}{x_{i_k}-\hat{x}_{i_{k-1}}}$, then $\hat{x}_{i_{k}}\not\in T_{P}$.
\end{claim}

\begin{proof}
Suppose by contradiction that  $\hat{x}_{i_{k}} \in T_{P}$. By Assumption \ref{ass: intersection_opt} - item 2, we must have that $\hat{x}_{i_{k}+1}\in \mathbb{X}^{\star}$ and that
%
\begin{equation}
\frac{\hat{v}_{i_{k}+1} - \hat{v}_{i_{k}}}{\hat{x}_{i_{k}+1} - \hat{x}_{i_{k}}} \geq \frac{Q^{q}(x_{i_{k}+1}) - Q^{m}(x_{i_{k}})}{\delta} = \frac{\hat{v}_{i_{k}+1} - \hat{v}_{i_k}}{\delta} > 
\frac{\hat{v}_{x_{i_k}} - \hat{v}_{x_{i_{k-1}}}}{x_{i_{k}} - x_{i_{k-1}}},
\end{equation}
%
where $q$ and $m$ are elements of $\DChSet$ that satisfy the notation from Definition \ref{def: TP} with $\hat{x}_{i_{k}}$ as the turning point. To complete the proof, since the above equation implies $i_{k}+1$ makes a left turn from the point $i_{k}$ and  $i_{k-1}$ on the value correspondence, by Line \ref{lst:line: left_turn} of Algorithm \ref{alg: NCE_scan}, we have that $\hat{x}_{i_{k}+1}\in \mathbb{X}$, $\hat{x}_{i_{k}+1}\in \mathbb{X}^{\star}$ and $i_{k+1} = i_{k}$. However, this contradicts the assumption of the claim that $x_{i_{k}}$ makes a right turn on the value correspondence. Thus, $\hat{x}_{i_{k}}\not\in T_{P}$. \end{proof}

The proof for the following claim is analogous to Claim \ref{clm: x_k_not_TP_1}, replacing $i_{k}+1$ with $l+1$. 

\begin{claim}\label{clm: jump_detect}
If $\hat{x}^{\prime}_{i_{l+1}} = Q^{\hat{\mVecVar}}(\hat{x}_{i_{l+1}})$, $\hat{\mVecVar} \not= \mVecVar$ and $\hat{x}^{\prime}_{i_{L+1}}\leq x^{\text{TP}}$, then $\hat{x}^{\prime}_{i_{l+1}}\not \in \mathbb{X}$. 
\end{claim}

\begin{proof}
Since $x^{\text{TP}}\geq \hat{x}_{i_{l+1}}$, we have:
%
% define sigma_m above
% define z_{j+1}
% why via the lipshitz condition? <- because sigma is the same between the two points 
\begin{equation*}
    V(\hat{x}_{i_k}) = G^{\mVecVar}(\hat{x}_{i_k}, \hat{z}_{i_k})
\end{equation*}
%
\begin{equation*}
    \hat{x}_{i_{l+1}}^{\prime} = f(\hat{x}_{i_{l+1}},\sigma^{\hat{\mVecVar}}(\hat{x}_{i_{l+1}}))
\end{equation*}

By Assumption \ref{ass: jump_size}, Assumption \ref{ass: lipshitz} and using the reverse triangle inequality, we have:
%
\begin{align*}
     \vert \hat{x}_{i_{l+1}}^{\prime} - \hat{x}_{i_k}^{\prime}\vert &  \geq \Big\vert \vert \hat{x}_{i_{l+1}}^{\prime} - f(\sigma^{\mVecVar}(\hat{x}_{i_{l+1}}),\hat{x}_{i_{l+1}})\vert  - \vert \hat{x}_{i_k}^{\prime}- f(\sigma^{\mVecVar}(\hat{x}_{i_{l+1}}), \hat{x}_{i_{l+1}})\vert \Big\vert \\ & \geq \vert D - M(\hat{x}_{i_{l+1}}- \hat{x}_{i_k})\vert 
\end{align*}
%
%
Dividing through, and noting Assumption \ref{ass: grid_big}, we get:
%
\begin{equation*}
    \frac{\vert \hat{x}_{i_{l+1}}^{\prime} - \hat{x}_{i_k}^{\prime}\vert}{\hat{x}_{i_{l+1}} - \hat{x}_{i_k}} \geq \vert \frac{D}{\hat{x}_{i_{l+1}}- \hat{x}_{i_k}}- M\vert >  \frac{D}{\delta} - M
\end{equation*}
%
However, since $\frac{D}{\delta}-M \geq \bar{M}$, by line \ref{lst:line: right_no_jump} of Algorithm \ref{alg: NCE_scan}, this implies $\hat{x}_{i_{l+1}}\not\in \mathbb{X}$. 
\end{proof}

%%%%%%%%%%%%%%%%%%%

\section{Additional tables and figures}

Figure \ref{fig: haus_smegm} in shows the upper envelope of the value correspondence recovered by FUES and the corresponding section of the housing policy function defined on total wealth at time $t$. 

\begin{figure} \centering \captionsetup{width=.7\linewidth} \hspace{-1.5cm}\includegraphics[scale = 0.75]{plots/hous_vf_aprime_all_small_59.png} \captionof{figure}{Value function and optimal housing policy function for the top right section of the endogenous grid in Figure \ref{fig: haus_bigegm}.} \label{fig: haus_smegm} \end{figure}

The following formalises the definition of the points on the space $S$ where choice-specific value functions intersect:
%
\begin{definition}\label{def: TP}
Let $T_{P}\subset S$ be the set of `crossing points' between the future choice-specific value functions. That is, the set of $x\in S$ such that for some $\mVecVar,\qVecVar\in \DChSet$, $Q^{\mVecVar}(x) = Q^{\qVecVar}(x)$ and for some $\epsilon$ with $\epsilon>0$, $V(x+\epsilon) = Q^{\qVecVar}(x+\epsilon)> Q^{\mVecVar}(x+\epsilon)$ and $V(x-\epsilon) = Q^{\mVecVar}(x-\epsilon)> Q^{\qVecVar}(x-\epsilon)$. 
\end{definition}

\vspace{0.5em}

\begin{table}[!ht]
\centering
\caption{\small Speed and accuracy of FUES in the continuous housing investment with adjustment frictions model with infinite horizon.}\label{table:housinginf}
\small
\begin{tabular}{ccccc|ccc}
\toprule
\multirow{2}{*}{\textit{Grid Size A}} & \multirow{2}{*}{\textit{$\tau$}} & \multicolumn{3}{c}{\textit{Time per iteration}} & \multicolumn{3}{c}{\textit{Euler error}} \\
& & \textbf{RFC} & \textbf{FUES} & \textbf{NEGM} & \textbf{RFC} & \textbf{FUES} & \textbf{NEGM} \\
\midrule
\multirow{3}{*}{300} & 0.03 & 7.713 & 6.871 & 11.336 & -3.799 & -3.506 & -2.338 \\
             & 0.07 & 7.639 & 6.734 & 10.932 & -3.037 & -2.945 & -2.061 \\
             & 0.15 & 7.846 & 6.829 & 10.833 & -2.357 & -2.303 & -1.680 \\
\midrule
\multirow{3}{*}{400} & 0.03 & 10.537 & 9.269 & 15.285 & -3.756 & -3.533 & -2.386 \\
             & 0.07 & 10.429 & 9.370 & 14.707 & -3.035 & -2.938 & -2.102 \\
             & 0.15 & 10.504 & 9.468 & 14.804 & -2.358 & -2.305 & -1.699 \\
\midrule
\multirow{3}{*}{500} & 0.03 & 12.290 & 10.609 & 17.247 & -3.727 & -3.408 & -2.402 \\
             & 0.07 & 11.998 & 10.558 & 16.975 & -3.028 & -2.862 & -2.106 \\
             & 0.15 & 12.044 & 10.589 & 20.991 & -2.356 & -2.298 & -1.704 \\
\midrule
\bottomrule
\end{tabular}

\vspace{1em}

\caption{\small Model parameters: $\gamma = 3$, $\alpha = 0.66$, $\beta = 0.93$, $\tau = 0.18$, $r = 0.01$, $\bar{\theta} = 1.34$. FUES jump detection threshold is $\bar{M} = 1.2$, with a 4-step forward and backward scan used. RFC jump detection threshold is $\bar{M} = 1.2$, with a search radius of $0.75$. Time per iteration is reported in seconds, while mean Euler error is reported in decimal logarithms \citep{JUDD1992410, Fella2014, Druedahl2017}.}
\end{table}

\vspace{-2em}

\begin{figure}[!htb]
\centering
\includegraphics[scale =0.7]{plots/ret_cons_all.png}
\captionsetup{width=.7\linewidth}
\captionof{figure}{Optimal consumption functions without smoothing for workers in Application 1.}\label{fig: finite_plot_pol}
\end{figure}



%%%%%%%%%%%%%%%%%%%%%%%%%%%%%%%%%%%%%%%%%%%%%

\end{document}